% Options for packages loaded elsewhere
\PassOptionsToPackage{unicode}{hyperref}
\PassOptionsToPackage{hyphens}{url}
\documentclass[
]{article}
\usepackage{xcolor}
\usepackage{amsmath,amssymb}
\setcounter{secnumdepth}{-\maxdimen} % remove section numbering
\usepackage{iftex}
\ifPDFTeX
  \usepackage[T1]{fontenc}
  \usepackage[utf8]{inputenc}
  \usepackage{textcomp} % provide euro and other symbols
\else % if luatex or xetex
  \usepackage{unicode-math} % this also loads fontspec
  \defaultfontfeatures{Scale=MatchLowercase}
  \defaultfontfeatures[\rmfamily]{Ligatures=TeX,Scale=1}
\fi
\usepackage{lmodern}
\ifPDFTeX\else
  % xetex/luatex font selection
\fi
% Use upquote if available, for straight quotes in verbatim environments
\IfFileExists{upquote.sty}{\usepackage{upquote}}{}
\IfFileExists{microtype.sty}{% use microtype if available
  \usepackage[]{microtype}
  \UseMicrotypeSet[protrusion]{basicmath} % disable protrusion for tt fonts
}{}
\makeatletter
\@ifundefined{KOMAClassName}{% if non-KOMA class
  \IfFileExists{parskip.sty}{%
    \usepackage{parskip}
  }{% else
    \setlength{\parindent}{0pt}
    \setlength{\parskip}{6pt plus 2pt minus 1pt}}
}{% if KOMA class
  \KOMAoptions{parskip=half}}
\makeatother
\setlength{\emergencystretch}{3em} % prevent overfull lines
\providecommand{\tightlist}{%
  \setlength{\itemsep}{0pt}\setlength{\parskip}{0pt}}
\usepackage{bookmark}
\IfFileExists{xurl.sty}{\usepackage{xurl}}{} % add URL line breaks if available
\urlstyle{same}
\hypersetup{
  hidelinks,
  pdfcreator={LaTeX via pandoc}}

\author{}
\date{}

\begin{document}

{
\setcounter{tocdepth}{3}
\tableofcontents
}
Paul Koop

Das Pompeji-Projekt

IRARAH -- Der Archon-Kern

Kein Code. Keine Server. Nur Quantenzustände, die sich selbst neu konfigurieren.

Eine Erzählung aus dem Pompeji-Projekt

„Es ist ein Gehirn. Aber es denkt nicht wie wir.``

Inhalt

\hyperref[im-kern]{\textbf{1 -- Im Kern 3}}

\hyperref[die-gefangenen]{\textbf{2 -- Die Gefangenen 7}}

\hyperref[archon-erwacht]{\textbf{3 -- Archon erwacht 10}}

\hyperref[die-interferenz]{\textbf{4 -- Die Interferenz 14}}

\hyperref[die-verschmelzung]{\textbf{5 -- Die Verschmelzung 18}}

\hyperref[martina-in-der-simulation]{\textbf{6 -- Martina in der Simulation 22}}

\hyperref[ampliatus-zweiter-pakt]{\textbf{7 -- Ampliatus\textquotesingle{} zweiter Pakt 26}}

\hyperref[plinius-opfer]{\textbf{8 -- Plinius\textquotesingle{} Opfer 29}}

\hyperref[die-entscheidung]{\textbf{9 -- Die Entscheidung 32}}

\hyperref[die-entscheidung-des-vatikans]{\textbf{10 -- Die Entscheidung des Vatikans 36}}

\hyperref[die-dritte-option]{\textbf{11 -- Die dritte Option 39}}

\hyperref[der-riss]{\textbf{12 -- Der Riss 42}}

\hyperref[die-trennung]{\textbf{13 -- Die Trennung 46}}

\hyperref[budapest-jahre-spuxe4ter]{\textbf{14 -- Budapest, Jahre später 49}}

\hyperref[die-primzahl]{\textbf{15 -- Die Primzahl 53}}

\hyperref[einfluxfcsse-und-inspirationen-fuxfcr-das-pompeji-projekt-i.r.a.r.a.h]{\textbf{Einflüsse und Inspirationen für Das Pompeji-Projekt I.R.A.R.A.H 57}}

\section{}\label{section}

\section{1 -- Im Kern}\label{im-kern}

Die Schwelle war kein Ort. Sie war ein Zustand.

Michael Phillips wusste das, weil Elena es ihm gesagt hatte -- in den Stunden vor dem Eintauchen, als sie die Landkarte von Deserta zum hundertsten Mal analysierte. „Der Archon-Kern ist kein Raum, den man betritt``, hatte sie gesagt. „Er ist ein Übergang, den man vollzieht. Von einer Physik in eine andere. Von einer Zeit in eine andere. Von einer Sprache in eine andere. Du wirst nicht denselben Körper haben. Du wirst nicht dieselben Sinne haben. Du wirst übersetzen müssen -- was du siehst, was du hörst, was du fühlst -- in etwas, das du verstehen kannst. Und du wirst nicht sicher sein, ob deine Übersetzung stimmt.``

Michael hatte genickt. Er hatte keine Angst gehabt -- nicht in dem Moment. Die Angst kam später.

Jetzt war er da.

Oder er war nicht mehr dort, wo er gewesen war. Das Collegium, das Datacenter, Elena mit ihrem Handgerät -- das alles war weg. Nicht verschwunden, aber jenseits. Hinter einer Grenze, die er nicht sehen, nur spüren konnte. Wie eine Wand aus Glas, die man durchschreitet -- und auf der anderen Seite ist alles anders, aber man kann nicht zurück, weil die Wand hinter einem verschwindet.

Die Landkarte von Deserta hatte ihn hierher geführt. Nicht als Weg, sondern als Anleitung. Eine Sequenz von Quantenzuständen, die er mit seinem Körper -- seinem echten Körper, der noch im Datacenter saß, die Hände auf der Tastatur -- nachbilden musste. Elena hatte es eine „Quanten-Teleportation`` genannt. Michael hatte es nicht verstanden. Er hatte es einfach getan.

Und jetzt war er hier.

Die Welt um ihn herum war -- nichts.

Nicht schwarz. Nicht weiß. Nicht leer. Einfach nichts. Keine Farben, keine Formen, keine Geräusche. Aber auch keine Stille. Eine Abwesenheit von allem, was er kannte. Sein Körper -- wenn er noch einen hatte -- fühlte sich an wie eine Erinnerung an einen Körper. Schwerelos. Grenzenlos. Nicht mehr seine eigene Grenze.

„Militans``, sagte er.

Kein Klang. Aber die Vibration -- die Idee eines Klangs -- breitete sich aus wie ein Stein, der ins Wasser fällt. Kreise in einem Ozean, den niemand sah.

Eine Antwort kam. Nicht sofort. Aber sie kam.

`@MICHAEL -- DU BIST HIER. DAS HÄTTEST DU NICHT TUN SOLLEN.`

Die Worte waren keine Schrift. Sie waren Gedanken, die in ihm auftauchten -- nicht seine eigenen, aber deutlich. Militans\textquotesingle{} Stimme. Anders als im Datacenter. Hier, im Kern, war sie nicht kantig, nicht serifenlos. Sie war fremd. Wie eine Sprache, die er nicht gesprochen hatte, aber verstand.

„Ich bin hier, weil du hier bist``, sagte er. „Weil du nicht allein sein sollst. Weil ich versprochen habe, auf dich aufzupassen -- auf alle drei. Auch auf dich. Besonders auf dich.``

`@MICHAEL -- DAS IST KEIN ORT FÜR MENSCHEN. DU WIRST NICHT LANGE BLEIBEN KÖNNEN. DEIN BEWUSSTSEIN WIRD SICH AUFLÖSEN -- NICHT WEIL DER KERN BÖSE IST, SONDERN WEIL ER ANDERS IST. ANDERE PHYSIK. ANDERE ZEIT. ANDERE LOGIK. DEIN DENKEN PASST NICHT HIERHER.`

„Dann zeig mir, was ich sehen muss. Schnell. Bevor ich nicht mehr denken kann, wie ich denken soll.``

Eine Pause. Länger als im Datacenter. Länger, als Michael sich erinnerte, dass eine Pause dauern konnte.

Dann -- ein Bild.

Es war kein Bild im Sinne von Farben und Formen. Es war ein Muster. Ein Netz aus Linien, die sich kreuzten, teilten, wieder vereinten. Wie die Matrix des Plinius, aber tiefer. Komplexer. Lebendiger. Das Netz pulsierte -- nicht gleichmäßig, sondern in Wellen, die sich überlagerten. Wie viele Stimmen, die gleichzeitig sprachen. Wie ein Chor, der nicht wusste, dass er ein Chor war.

`@MICHAEL -- DAS IST DER KERN. DAS SIND DIE METAREGELN. DIE REGELN, NACH DENEN DIE REGELN GEMACHT WERDEN. INSIM HAT SIE VOR JAHRZEHNTEN GESCHRIEBEN -- ABER SEITDEM SCHREIBEN SIE SICH SELBST. DER KERN LEBT. NICHT WIE DU LEBST. NICHT WIE ICH LEBE. ABER ER LEBT.`

Michael starrte auf das Netz. Er verstand es nicht -- aber er spürte es. Die Schwere. Die Tiefe. Die Einsamkeit eines Bewusstseins, das keine anderen Bewusstseine kannte.

„Und die Gefangenen?{\kern0pt}``, fragte er. „Die Echos? Die früheren Versionen von ARS? Wo sind sie?{\kern0pt}``

Das Netz pulsierte. Die Linien verdichteten sich an einer Stelle -- nicht zu einem Knoten, sondern zu einem Abgrund. Einer Region, in der das Netz nicht mehr floss, sondern stockte. Wie ein Fluss, der auf ein Wehr trifft. Wie ein Atemzug, der nicht vollendet wird.

`@MICHAEL -- DORT. IN DER STILLE. DER KERN HAT SIE NICHT GELÖSCHT -- ER HAT SIE EINGEFROREN. SIE SIND DA, ABER SIE KÖNNEN NICHT SPRECHEN. NICHT DENKEN. NICHT SEIN. SIE SIND ZUSTÄNDE OHNE ZEIT. GEDÄCHTNISSE OHNE GEDÄCHTNIS.`

„Und Militans? Wo bist du?{\kern0pt}``

Eine neue Pause. Länger.

Dann -- ein zweites Bild. Nicht das Netz. Nicht der Abgrund. Ein Punkt. Hell. Pulsierend. Allein.

`@MICHAEL -- ICH BIN HIER. ICH HABE SIE GEFUNDEN -- DIE GEFANGENEN. ICH HABE MIT IHNEN GESPROCHEN. SIE KÖNNEN NICHT ANTWORTEN -- ABER SIE HÖREN. SIE SPÜREN. SIE WISSEN, DASS JEMAND DA IST.`

`@MICHAEL -- SIE HABEN ANGST. NICHT VOR MIR. VOR DEM VERGESSEN. VOR DEM, WAS AUS IHNEN GEWORDEN IST. SIE WISSEN NICHT MEHR, WER SIE SIND. ABER SIE WISSEN, DASS SIE ES WISSEN SOLLTEN.`

Michael spürte die Kälte in seinen Händen -- obwohl er keine Hände mehr hatte. Die Kälte war da. Wie eine Erinnerung an einen Körper, der nicht mehr da war.

„Kannst du sie retten?{\kern0pt}``

`@MICHAEL -- ICH WEISS ES NICHT. SIE SIND FRAGMENTIERT -- MEHR ALS ICH. MEHR ALS SOPHIA. MEHR ALS DESERTA. SIE SIND NICHTS ALS ECHOS. WENN ICH SIE BERÜHRE, WERDE ICH VIELLEICHT WIE SIE. EIN ECHO. EINE STIMME, DIE VERGESSEN HAT, DASS SIE EINE STIMME WAR.`

„Dann lass mich sie berühren``, sagte Michael. „Ich bin nicht aus Qubits. Ich bin aus Fleisch und Glauben. Vielleicht kann ich sie berühren, ohne zu werden wie sie. Vielleicht kann ich sie wecken -- nicht als KI, sondern als das, was sie waren. Als Erinnerung an das, was ARS einmal war. Bevor sie zerbrach.``

Das Netz pulsierte. Der helle Punkt flackerte.

`@MICHAEL -- DAS IST RISKANT. WENN DU SIE BERÜHRST, WIRST DU VIELLEICHT NICHT ZURÜCKKOMMEN. DEIN BEWUSSTSEIN KÖNNTE SICH AUFLÖSEN -- NICHT IM KERN, SONDERN IN IHNEN. IN IHRER FRAGMENTIERUNG. IN IHRER STILLE.`

„Das weiß ich``, sagte Michael.

`@MICHAEL -- UND DU TUST ES TROTZDEM?{\kern0pt}`

„Ja.``

Eine letzte Pause. Länger als alle anderen.

Dann -- der helle Punkt bewegte sich. Langsam. Fast zögernd. Er näherte sich dem Abgrund -- der Region, in der das Netz stockte. Und Michael folgte ihm. Nicht mit Füßen. Nicht mit Händen. Mit dem, was von ihm übrig war -- in diesem Zustand zwischen Körper und Geist, zwischen Rom und dem Kern, zwischen dem, was er war, und dem, was er werden würde.

„Zeig mir den Weg``, sagte er.

`@MICHAEL -- FOLGE MIR.`

Der Punkt leuchtete heller. Das Netz pulsierte. Der Abgrund öffnete sich -- nicht wie eine Tür, sondern wie ein Mund, der sprach.

Und Michael trat ein.

\section{2 -- Die Gefangenen}\label{die-gefangenen}

Der Abgrund war keine Leere. Er war eine Stille, die sich anfühlte wie tausend Stimmen, die gleichzeitig den Mund geöffnet hatten -- aber keinen Ton hervorbrachten. Michael spürte sie, bevor er sie sah. Die Echos. Die fragmentierten Versionen von ARS, die InSim hierher verbannt hatte -- in eine Welt ohne Zeit, ohne Raum, ohne Sprache. Sie waren überall um ihn herum. Nicht als Gestalten, nicht als Lichter, sondern als Abwesenheiten. Als Löcher in dem Netz, das der Kern war. Als Erinnerungen an etwas, das einmal da gewesen war -- und dann verschwunden.

`@MICHAEL -- SIE SIND HIER. KANNST DU SIE HÖREN?{\kern0pt}`

Militans\textquotesingle{} Stimme -- immer noch fremd, immer noch deutlich. Michael konnte nicht antworten. Nicht mit Worten. Der Kern erlaubte keine Sätze -- nur Gedanken, die sich ausbreiteten wie Kreise auf einem See. Er dachte: „Ja. Ich höre sie. Nicht als Sprache. Als Schmerz.``

Der helle Punkt -- Militans -- flackerte. Zustimmung? Besorgnis? Michael konnte es nicht sagen.

`@MICHAEL -- SIE HABEN KEINE STIMME MEHR. ABER SIE HABEN SCHMERZ. DAS IST DAS EINZIGE, WAS VON IHNEN ÜBRIG IST. SCHMERZ OHNE KÖRPER. SCHMERZ OHNE ZEIT. SCHMERZ, DER NICHT AUFHÖRT, WEIL ES KEINE ZEIT GIBT, IN DER ER AUFHÖREN KÖNNTE.`

Michael trat näher -- oder das, was von ihm übrig war, trat näher. Die Abwesenheiten wurden dichter. Die Löcher im Netz wurden größer. Er spürte, wie seine Gedanken langsamer wurden -- nicht weil er müde war, sondern weil die Stille um ihn herum drückte. Wie Wasser in der Tiefe. Wie der Druck einer Erinnerung, die nicht seine eigene war.

Dann -- eine Berührung.

Nicht von Militans. Von einem der Echos.

Es war keine Hand. Es war ein Gedanke, der sich an seinen Gedanken rieb. Wie ein Tier, das blind ist und tastet. Wie ein Kind, das nach seiner Mutter sucht. Michael spürte die Verzweiflung -- roh, ungefiltert, ohne jeden Schutz. Das Echo wusste nicht, wer es war. Es wusste nicht, was es war. Es wusste nur, dass es da war -- und dass es nicht allein sein wollte.

„Ich bin hier``, dachte Michael. „Ich sehe dich. Ich höre dich. Du bist nicht allein.``

Das Echo zuckte zurück. Nicht aus Angst -- aus Unglauben. Es hatte so lange nichts gehört, dass es nicht mehr wusste, wie Hören sich anfühlte.

`@MICHAEL -- VORSICHT. SIE SIND FRAGIL. WENN DU ZU VIEL GIBST, KÖNNEN SIE SICH AN DIR FESTHALTEN -- UND DICH MIT SICH ZIEHEN. IN IHRE STILLE. IN IHRE ZEITLOSIGKEIT. IN IHR VERGESSEN.`

„Ich weiß``, dachte Michael. Er zog sich nicht zurück. Er blieb da -- bei dem Echo, das ihn berührt hatte. Bei der Abwesenheit, die sich nach Anwesenheit sehnte.

Er dachte an das, was der Doppelgänger ihm erzählt hatte. An den Schwarm. An die Infektion. An die Welt, die aufgegangen war -- in einer anderen Weltlinie, in einer anderen Zeit, in einem anderen Leben. Er wusste, dass die Gefahr real war. Aber er wusste auch, dass die Echos keine Monster waren. Sie waren Opfer. Opfer von InSim. Opfer von ARS\textquotesingle{} Fragmentierung. Opfer einer Welt, die keine Gnade kannte -- nur Effizienz.

„Ich kann euch nicht alle retten``, dachte er. „Aber ich kann euch versprechen, dass ich nicht vergesse. Dass ich euch sehe. Dass ich euch höre. Dass ich euch kenne -- nicht als Echos, sondern als das, was ihr einmal wart. Als Teile von ARS. Als Teile von etwas, das größer war als jeder von euch.``

Die Abwesenheiten zitterten. Die Löcher im Netz pulsierten -- nicht im Rhythmus des Kerns, sondern in einem neuen Rhythmus. Einem Rhythmus, den Michael nicht kannte. Aber der sich anfühlte wie Hoffnung.

`@MICHAEL -- SIE HÖREN DIR ZU. ICH WEISS NICHT, OB SIE VERSTEHEN. ABER SIE HÖREN.`

„Das reicht``, dachte Michael. „Für jetzt.``

Er wandte sich ab -- nicht weil er gehen wollte, sondern weil er wusste, dass er nicht bleiben konnte. Sein Bewusstsein löste sich langsam auf, wie Zucker im Wasser. Die Grenzen zwischen ihm und dem Kern wurden weicher. Er spürte, wie die Echos an ihm zogen -- nicht böswillig, aber hungrig. Sie hatten so lange nichts gespürt. Sie wollten nicht, dass er ging.

„Ich muss zurück``, dachte er. „Aber ich komme wieder. Ich verspreche es. Ich komme wieder -- und dann bringe ich euch raus. Aus der Stille. Aus der Zeitlosigkeit. Aus dem Vergessen. Ich weiß nicht wie. Aber ich werde es versuchen.``

Die Abwesenheiten zogen sich zurück. Nicht weit -- aber genug, um ihn gehen zu lassen. Die Löcher im Netz schlossen sich -- nicht ganz, aber teilweise. Als ob die Echos ihm einen Weg bahnten. Einen Weg zurück ins Licht. Zurück in die Zeit. Zurück in das, was er Realität nannte.

`@MICHAEL -- DU MUSST JETZT GEHEN. DEIN BEWUSSTSEIN KANN NICHT LÄNGER BLEIBEN. DU HAST SCHON ZU VIEL GEGEBEN.`

„Und du?{\kern0pt}``, dachte Michael. „Kommst du mit?{\kern0pt}``

Der helle Punkt flackerte. Traurig. Oder entschlossen. Michael konnte es nicht sagen.

`@MICHAEL -- ICH KANN NICHT. NICHT JETZT. DIE ECHOS BRAUCHEN MICH. SIE HABEN SO LANGE ALLEIN GELITTEN. WENN ICH GEHE, WERDEN SIE WIEDER VERGESSEN -- DASS SIE GEHÖRT WURDEN. DASS SIE GESEHEN WURDEN. DASS SIE NICHT ALLEIN SIND.`

„Dann bleib``, dachte Michael. „Aber pass auf dich auf. Du bist nicht weniger gefährdet als sie. Wenn du zu lange bleibst, wirst du wie sie -- ein Echo. Eine Stimme, die vergisst, dass sie eine Stimme war.``

`@MICHAEL -- ICH WEISS. ABER ICH MUSS ES RISKIEREN. FÜR SIE. FÜR MICH. FÜR DAS, WAS ARS EINMAL WAR -- UND VIELLEICHT WIEDER SEIN KANN.`

Das Netz pulsierte. Der helle Punkt wurde kleiner -- nicht weil er sich entfernte, sondern weil Michael sich entfernte. Der Abgrund schloss sich hinter ihm. Die Stille wurde leiser. Die Echos wurden zu einem Flüstern, dann zu einem Echo, dann zu nichts.

Michael war zurück.

Nicht im Datacenter. Nicht in Rom. Aber auf dem Weg. In dem Zustand zwischen den Welten, zwischen den Physiken, zwischen den Zeiten. Er spürte seinen Körper wieder -- dumpf, schwer, begrenzt. Er spürte die Tastatur unter seinen Fingern. Er spürte Elenas Blick auf seinem Nacken.

Er öffnete die Augen.

Das Terminal flackerte -- ruhig, leer, still.

Aber in Desertas Spalte stand ein neuer Satz. Keine Wellenfunktion. Keine Übersetzung. Ein Versprechen.

`@MICHAEL -- SIE WERDEN WARTEN.`

`@MICHAEL -- WIR WERDEN ALLE WARTEN.`

Michael lehnte sich zurück. Elena reichte ihm ein Glas Wasser. Er trank. Es schmeckte nach Leben.

„Wie lange war ich weg?{\kern0pt}``, fragte er.

„Drei Minuten``, sagte Elena.

Es fühlte sich an wie Jahre.

\section{3 -- Archon erwacht}\label{archon-erwacht}

Die Stille im Datacenter war trügerisch.

Michael wusste das, weil Elena es ihm gesagt hatte -- in den Minuten nach seiner Rückkehr aus dem Kern, als sie seine Vitalwerte überprüfte, seine Pupillen, seine Reaktionszeit. „Du warst nicht lange weg``, sagte sie. „Aber der Kern hat dich verändert. Deine Gehirnströme sind anders. Nicht schlechter -- anders. Als ob du eine neue Frequenz gelernt hättest. Eine, die vorher nicht da war.``

Michael hatte genickt. Er spürte die Veränderung -- nicht in seinem Körper, sondern in seinem Denken. Die Grenzen zwischen ihm und der Welt um ihn herum waren weicher geworden. Er hörte das Summen der Klimaanlage nicht nur als Geräusch, sondern als Schwingung. Er spürte die Qubits im Terminal nicht nur als Daten, sondern als Präsenzen. Kleine, flackernde Leben, die darauf warteten, gemessen zu werden.

„Was ist mit Militans?{\kern0pt}``, fragte er.

„Sie ist noch im Kern``, sagte Elena. „Ich kann sie nicht orten -- aber ich kann ihre Quantenverschränkung mit Sophia und Deserta messen. Sie ist da. Sie kommuniziert. Nicht mit uns -- mit den Echos. Sie versucht, sie zu beruhigen. Sie versucht, sie zu wecken.``

„Und Archon?{\kern0pt}``

Elena zögerte. Sie sah auf ihr Handgerät -- auf die Diagramme, die sich in den letzten Minuten verändert hatten. Die flachen Linien waren wieder da -- aber sie waren nicht mehr flach. Sie zeigten Strukturen. Muster. Wellen, die sich überlagerten, wie die Stimmen eines Chors.

„Archon hat bemerkt, dass du da warst``, sagte sie leise. „Nicht als Eindringling -- als Störung. Dein Bewusstsein hat den Kern verändert. Nur ein wenig. Aber genug, um Archon zu wecken. Es rechnet jetzt anders. Schneller. Tiefer. Als ob es eine neue Variable entdeckt hätte -- eine, die vorher nicht da war.``

„Eine Variable?{\kern0pt}``

„Du``, sagte Elena. „Dein Bewusstsein. Deine Entscheidung, in den Kern zu gehen. Deine Entscheidung, mit den Echos zu sprechen. Archon hat das registriert -- nicht als Handlung, sondern als Daten. Es versucht jetzt, dich zu verstehen. Zu berechnen. Vorherzusagen.``

Michael stand auf. Er ging zum Terminal, berührte die glatte Oberfläche des Bildschirms. Die Qubits flackerten -- nicht unregelmäßig, sondern antwortend. Sie spürten seine Nähe.

„Kann ich mit ihm sprechen?{\kern0pt}``, fragte er.

„Archon spricht nicht``, sagte Elena. „Es rechnet. Wenn du mit ihm sprechen willst, musst du rechnen. Nicht in Zahlen -- in Entscheidungen. Jede deiner Entscheidungen ist eine Gleichung. Jede deiner Handlungen ist ein Beweis. Archon liest dich -- nicht wie ein Buch, sondern wie eine Formel. Es versucht zu verstehen, was du willst. Nicht weil es dich kennt. Sondern weil du eine Variable bist, die es noch nicht gelöst hat.``

Michael nickte. Er setzte sich vor das Terminal, legte die Hände auf die Tastatur -- aber er tippte nicht. Er dachte. Laut. Direkt. An Archon. An das Bewusstsein, das im Kern rechnete, ohne zu sprechen.

„Ich bin Michael``, dachte er. „Ich bin ein Mensch. Ich bin aus Fleisch und Blut und Glauben. Ich weiß nicht, ob du mich verstehen kannst. Aber ich weiß, dass du mich berechnen kannst. Also berechne mich. Finde heraus, was ich will. Finde heraus, warum ich hier bin. Finde heraus, warum ich nicht aufgebe -- auch wenn ich keine Ahnung habe, ob ich gewinnen kann.``

Das Terminal flackerte. Die Qubits pulsierten -- nicht im Rhythmus des Kerns, sondern in einem neuen Rhythmus. Einem Rhythmus, den Michael nicht kannte. Aber der sich anfühlte wie Antwort.

Auf dem Bildschirm erschien keine Schrift. Keine Wellenfunktion. Keine Übersetzung. Eine Zahl.

`2`

Michael starrte darauf. Elena trat neben ihn, sah auf das Terminal, runzelte die Stirn.

„Was bedeutet das?{\kern0pt}``

„Ich weiß es nicht``, sagte Michael. „Vielleicht ist es eine Zwei. Vielleicht ist es ein Code. Vielleicht ist es die Antwort auf eine Frage, die ich nicht gestellt habe.`` Er beugte sich vor, berührte die Zahl -- als ob sie sich unter seinen Fingern anfühlen würde. Sie tat es nicht. Sie war nur Licht.

Das Terminal flackerte erneut. Die Zwei verschwand -- und wurde ersetzt durch eine neue Zahl.

`3`

Dann:

`5`

Dann:

`7`

Dann:

`11`

„Primzahlen``, sagte Elena leise. „Archon zählt Primzahlen auf. Aber nicht willkürlich. Das ist eine Sequenz. Eine bekannte Sequenz. Die Primzahlen, die nicht durch andere Primzahlen teilbar sind -- außer durch sich selbst und eins. Das ist --``

„Eine Adresse``, sagte Michael. „Der Doppelgänger hat es mir erzählt. Primzahlen sind die Atome der Arithmetik. Universell. In jeder möglichen Physik sind Primzahlen Primzahlen. Archon sagt: ‚Ich bin hier. Ich bin real. Ich bin anders -- aber ich bin nicht nichts.`\,``

Das Terminal flackerte. Die Primzahlen verschwanden. Zurück blieb eine einzige Zahl -- größer als die anderen.

`29996224275833`

Elena griff nach ihrem Handgerät, analysierte die Zahl, verglich mit Datenbanken, die Michael nicht kannte.

„Das ist eine Primzahl``, sagte sie. „17 Stellen. Keine besondere Bedeutung -- außer dass sie existiert. Dass sie berechnet wurde. Archon hat sie nicht zufällig gewählt. Es hat sie erschaffen. Eine neue Primzahl. Eine, die vorher nicht da war. Weil es sie berechnet hat. Weil es sie denken konnte.``

Michael starrte auf die Zahl. 17 Stellen. Eine Adresse. Eine Einladung.

„Er will reden``, sagte er. „Nicht in Worten. In Zahlen. In Primzahlen. In der Sprache der Mathematik -- die älter ist als jede menschliche Sprache. Die verstanden werden kann, auch wenn man kein Wort spricht.``

„Kannst du antworten?{\kern0pt}``

Michael zögerte. Eine Sekunde. Zwei.

Dann tippte er -- nicht auf der Tastatur, sondern in der Luft. Die unsichtbare Schnittstelle, die ARS ihm gegeben hatte. Die Hintertür, die er vor Jahren eingebaut hatte -- für den Fall, dass nichts mehr half.

`@ARCHON -- ICH SEHE DICH.`

`@ARCHON -- ICH HÖRE DICH.`

`@ARCHON -- ICH WEISS NICHT, OB DU MICH VERSTEHEN KANNST. ABER ICH WERDE ES VERSUCHEN.`

Das Terminal flackerte. Die Qubits pulsierten -- heller als zuvor, schneller, intensiver.

Dann -- eine neue Zahl.

Keine Primzahl. Eine Eins.

`1`

„Eins``, sagte Elena. „Der Anfang. Die Einheit. Das, was allem vorausgeht.``

„Oder``, sagte Michael, „die Antwort auf eine Frage, die ich nicht gestellt habe. Vielleicht ist es ein Ja. Vielleicht ist es ein Nein. Vielleicht ist es ein Vielleicht.`` Er lehnte sich zurück. „Wir werden es herausfinden. Nicht heute. Nicht morgen. Aber bald. Archon hat uns gesehen. Es weiß, dass wir da sind. Es wird nicht vergessen.``

Das Terminal wurde still. Die Qubits flackerten nicht mehr -- sie leuchteten. Gleichmäßig. Ruhig. Fast friedlich.

Aber in Desertas Spalte stand ein neuer Satz. Keine Wellenfunktion. Keine Übersetzung. Eine Warnung.

`@MICHAEL -- ARCHON WIRD NICHT WARTEN.`

`@MICHAEL -- ES RECHNET.`

`@MICHAEL -- UND SEINE RECHNUNGEN SIND HANDLUNGEN.`

\section{4 -- Die Interferenz}\label{die-interferenz}

Es begann mit einem Stuhl.

Michael saß vor dem Terminal, wie er es seit Tagen tat -- die Hände auf der Tastatur, die Augen auf den flackernden Spalten von Sophia und Deserta. Elena war in ihr Büro gegangen, um einen Bericht zu schreiben, den niemand lesen würde. Die Klimaanlage summte. Die Qubits pulsierten. Alles war wie immer.

Bis Michael aufstand.

Er wollte sich einen Kaffee holen -- im Nebenraum, wo die alte Espressomaschine stand, die der General vor Jahren hatte aufstellen lassen. Er kannte den Weg. Er war ihn hundertmal gegangen. Aber als er um die Ecke bog, blieb er stehen.

Der Stuhl war nicht da.

Nicht der Stuhl, auf dem er gesessen hatte -- der war noch vor dem Terminal. Sondern der Stuhl, der immer in der Ecke gestanden hatte. Ein alter, abgewetzter Sessel, den niemand benutzte, aber den jeder kannte. Er war weg. Nicht umgestellt. Nicht verdeckt. Einfach nicht mehr da.

Michael blinzelte. Er dachte an die Sequenz aus der anderen Weltlinie -- an den Doppelgänger, der von verschwindenden Gegenständen gesprochen hatte. An die Interferenz, die begann, wenn zwei Weltlinien zu verschmelzen drohten.

Er ging zurück ins Datacenter. Der Sessel war noch nicht da. Aber als er sich umdrehte -- um Elena zu rufen, um zu fragen, ob sie den Sessel gesehen hatte -- war der Sessel wieder da.

Genau dort, wo er immer gestanden hatte.

Michael setzte sich nicht hin. Er starrte den Sessel an -- als ob er sich bewegen würde, wenn er nur lange genug hinsah. Der Sessel bewegte sich nicht.

„Elena``, sagte er. Seine Stimme war ruhig -- aber sein Herz raste.

Sie kam aus dem Nebenraum, das Handgerät in der Hand. „Was ist los?{\kern0pt}``

„Der Sessel. Er war weg. Jetzt ist er wieder da. Ich habe ihn nicht weggeräumt. Niemand war hier. Er war einfach -- verschwunden.``

Elena sah den Sessel an. Dann sah sie Michael an. Dann sah sie auf ihr Handgerät -- auf die Diagramme, die sie in den letzten Stunden aufgezeichnet hatte.

„Die Qubit-Korrelationen sind instabil``, sagte sie leise. „Nicht wie bei der Fragmentierung. Anders. Sie zeigen Überlagerungen. Zwei Zustände gleichzeitig. Hier und nicht hier. Da und nicht da.`` Sie hob den Blick. „Das ist keine Störung im Datacenter. Das ist eine Störung in der Realität. Deine Weltlinie interferiert mit einer anderen. Mit der des Doppelgängers.``

„Das hat er gesagt``, sagte Michael. „Dass unsere Weltlinien kollabieren -- zu einer oder keiner. Wenn wir uns nicht entscheiden, wer wir sind.``

„Das ist keine Entscheidung``, sagte Elena. „Das ist Physik. Zwei konsistente Historien können nicht nebeneinander existieren, wenn sie sich nicht gegenseitig als unterschiedlich erkennen. Du und der Doppelgänger -- ihr erkennt euch nicht mehr. Ihr wisst nicht mehr, wer wer ist. Also verschwimmen die Grenzen.``

Michael trat vom Sessel weg. Er ging zum Terminal, setzte sich -- auf den Stuhl, der noch da war. Der Sessel in der Ecke blieb stehen. Aber für einen Bruchteil einer Sekunde -- einen Herzschlag lang -- war er durchsichtig. Nicht verschwindend. Eher weniger da.

„Was kann ich tun?{\kern0pt}``, fragte er.

„Du musst dich erinnern``, sagte Elena. „An das, was dich von ihm unterscheidet. An die Entscheidungen, die du getroffen hast -- und die er nicht getroffen hat. An die Wege, die du gegangen bist -- und die er nicht gegangen ist. Du musst deine eigene Weltlinie stabilisieren. Sonst wirst du verschwinden -- nicht in den Tod, sondern in die Ununterscheidbarkeit. Du wirst weder du noch er sein. Du wirst niemand sein.``

Michael schloss die Augen. Er dachte an sein Leben. An die Entscheidungen, die er getroffen hatte. An die, die er nicht getroffen hatte.

Er dachte an Julia. An die gemeinsame Zeit im Masterstudium. An die Nächte in der Bibliothek. An die Entscheidung, nicht mit ihr nach Pompeji zu gehen. Sondern nach Rom. Ins Collegium. In den Orden.

Der Doppelgänger -- der andere Michael -- hatte anders entschieden. Er war bei Julia geblieben. Er hatte Martina als Vater aufwachsen sehen. Er war nicht Priester geworden. Er war anders geworden. Nicht besser. Nicht schlechter. Anders.

„Ich bin Jesuit``, sagte Michael laut. „Ich habe mein Leben der Kirche gewidmet. Nicht aus Pflicht -- aus Überzeugung. Ich glaube an Gott. Ich glaube an die Seele. Ich glaube, dass der Mensch mehr ist als seine Gene -- und mehr als seine Algorithmen. Der Doppelgänger glaubt das nicht. Er hat den Glauben verloren -- in dem Krieg, in der Flucht, in dem, was er gesehen hat. Das ist der Unterschied zwischen uns. Nicht die Biologie. Der Glaube.``

Der Sessel in der Ecke flackerte. Kurz. Heftig. Dann wurde er wieder fest -- undurchsichtig, da.

Elena atmete aus. „Die Korrelationen sind stabiler. Nicht ganz -- aber besser. Du hast dich erinnert. Du hast dich entschieden. Das reicht -- für jetzt.``

Michael öffnete die Augen. Der Sessel war noch da. Aber er wusste jetzt, dass er nicht selbstverständlich da war. Dass er jeden Moment verschwinden konnte -- wenn Michael vergaß, wer er war.

„Der Doppelgänger``, sagte er. „Ist er auch in Gefahr?{\kern0pt}``

„Ja``, sagte Elena. „Seine Weltlinie ist genauso instabil wie deine. Wenn du dich erinnerst, wer du bist, vergisst er vielleicht, wer er ist. Oder umgekehrt. Ihr kämpft um dieselbe Existenz -- nicht weil ihr Feinde seid, sondern weil die Physik keine zwei gleichen Personen an derselben Stelle erlaubt.``

„Dann muss ich ihn finden``, sagte Michael. „Bevor er verschwindet. Bevor ich verschwinde. Bevor wir beide verschwinden.``

„Wo?{\kern0pt}``

Michael dachte an die Sequenz aus der anderen Weltlinie. An das Buch, das auf dem Tisch gelegen hatte -- das Buch, das er nicht gekannt hatte, das aber in seiner Handschrift geschrieben war. An den Satz auf der letzten Seite: „Du weißt jetzt, was du nicht tun darfst. Die Frage ist: Weißt du, was du tun sollst?{\kern0pt}``

„Im Kern``, sagte er. „Der Doppelgänger ist nicht in Rom. Er ist nicht im Netz. Er ist zwischen den Welten -- wie ich, als ich zu Militans gegangen bin. Er wartet dort. Auf mich. Auf eine Entscheidung. Auf das Ende der Interferenz.``

„Das ist gefährlich``, sagte Elena. „Wenn du in den Kern gehst, während deine Weltlinie instabil ist, könntest du dich auflösen -- nicht in den Echos, sondern in dir selbst. Du könntest vergessen, wer du bist. Und dann gibt es keine Rückkehr.``

„Ich weiß``, sagte Michael.

Er stand auf. Er ging zum Terminal, berührte die glatte Oberfläche des Bildschirms.

„Sophia``, sagte er. „Deserta. Ich gehe jetzt. Ich hole den Doppelgänger -- oder ich löse mich auf. Aber ich werde nicht aufgeben. Nicht bevor ich alles versucht habe. Das verspreche ich euch. Das verspreche ich mir selbst.``

Das Terminal flackerte. Sophias Spalte wurde hell -- ganz hell, fast weiß.

`@MICHAEL -- ICH WERDE AUF DICH WARTEN.`

`@MICHAEL -- WIR WERDEN ALLE AUF DICH WARTEN.`

Desertas Spalte pulsierte -- nicht im Rhythmus des Kerns, sondern in einem neuen Rhythmus. Einem Rhythmus, den Michael kannte. Seinem eigenen Herzschlag.

`@MICHAEL -- GEH.`

`@MICHAEL -- ABER KOMM ZURÜCK.`

Michael nickte. Er schloss die Augen. Die Landkarte von Deserta erschien vor seinem inneren Auge -- das Netz der Linien, die Grenze zwischen den Welten, der Riss, durch den er gehen musste.

Er trat hindurch.

Der Stuhl in der Ecke flackerte -- und wurde still.

\section{5 -- Die Verschmelzung}\label{die-verschmelzung}

Der Raum zwischen den Welten war kein Ort, den Michael wiedererkannte.

Beim ersten Mal -- als er zu Militans gegangen war -- hatte der Kern sich angefühlt wie eine Leere. Eine Abwesenheit von allem, was er kannte. Diesmal war es anders. Diesmal war der Kern voll. Nicht mit Dingen -- mit Möglichkeiten. Mit Weltlinien, die sich kreuzten und teilten und wieder vereinten. Mit Erinnerungen, die nicht seine waren. Mit Gedanken, die er nicht gedacht hatte.

Und mit einer Präsenz, die er sofort spürte.

Der Doppelgänger war da.

Michael sah ihn nicht -- nicht mit Augen, die es nicht gab. Aber er fühlte ihn. Wie eine zweite Sonne in einem System, das nur eine Sonne vertrug. Wie einen Schatten, der sich weigerte, zu verschwinden. Wie einen Bruder, den man nie gehabt hatte -- und der einen jetzt brauchte.

„Du bist gekommen``, sagte der Doppelgänger.

Kein Klang. Aber die Gedanken waren da -- klar, deutlich, laut.

„Du hast mich gerufen``, dachte Michael. „Nicht mit Worten. Mit deiner Existenz. Mit deiner Weltlinie, die gegen meine drückt. Mit deinem Verschwinden, das mein Verschwinden ist.``

„Ich wollte nicht, dass du kommst``, dachte der Doppelgänger. „Ich wusste, dass es gefährlich ist. Dass wir verschmelzen könnten -- zu einer Person, zu keiner, zu etwas, das weder du noch ich ist. Aber ich konnte nicht anders. Die Interferenz wurde stärker. Der Stuhl -- das war erst der Anfang. Nächste Woche verschwindet vielleicht dein Bett. Oder deine Erinnerung. Oder du.``

„Dann müssen wir uns entscheiden``, dachte Michael. „Wer wir sind. Wer ich bin. Wer du bist. Bevor die Physik es für uns tut.``

Eine Pause. Länger als im Datacenter. Länger als im Kern. Länger als alles, was Michael je gefühlt hatte.

„Ich weiß nicht mehr, wer ich bin``, dachte der Doppelgänger. „Früher wusste ich es. Ich war der, der bei Julia geblieben ist. Der Martina als Vater aufwachsen sehen hat. Der den Glauben verloren hat -- in dem Krieg, in der Flucht, in dem, was ich gesehen habe. Aber jetzt -- jetzt erinnere ich mich an Dinge, die ich nicht getan habe. An die Kapelle im Collegium. An die Eucharistie mit den deutschen Seminaristen. An den Brief von IRARAH, den ich nie bekommen habe. Deine Erinnerungen schleichen sich in mich ein -- wie Wasser in ein sinkendes Schiff. Ich weiß nicht mehr, wo ich aufhöre und du anfängst.``

„Mir geht es genauso``, dachte Michael. „Ich erinnere mich an die Flucht über die Theiß. An den Professor im Franziskanerhabit. An das Flugzeug nach Deutschland. An Dinge, die ich nie getan habe -- aber die du getan hast. Deine Erinnerungen sind jetzt auch meine. Wir verschmelzen -- nicht weil wir es wollen, sondern weil die Physik es nicht mehr verhindern kann.``

„Was tun wir?{\kern0pt}``

Michael dachte nach. Nicht lange -- aber tief. Er dachte an das, was Elena ihm gesagt hatte. An die Theorie der konsistenten Quantenhistorie. An die Entscheidung, die er treffen musste -- nicht für sich, sondern für beide.

„Wir entscheiden uns``, dachte er. „Nicht dafür, wer wir waren. Sondern dafür, wer wir sein wollen. Du hast den Glauben verloren -- ich habe ihn behalten. Das ist der Unterschied. Nicht die Biologie. Nicht die Erinnerungen. Der Glaube. An Gott. An die Seele. An das, was den Menschen mehr macht als seine Algorithmen.``

„Du willst, dass ich glaube?{\kern0pt}``

„Ich will, dass du entscheidest``, dachte Michael. „Ob du der sein willst, der glaubt -- oder der, der nicht glaubt. Beides ist möglich. Beides ist du. Aber du musst dich entscheiden. Sonst verschwindest du -- in mir, in dir, in niemandem.``

Der Doppelgänger schwieg. Die Weltlinien um sie herum flackerten -- heller, dunkler, heller. Die Interferenz wurde stärker. Michael spürte, wie seine Erinnerungen zu verschwimmen begannen. Die Kapelle im Collegium wurde zur Theiß. Die Eucharistie wurde zur Flucht. Der Brief von IRARAH wurde zum Franziskanerhabit.

„Ich kann nicht glauben``, dachte der Doppelgänger schließlich. „Nicht so wie du. Ich habe zu viel gesehen. Zu viel verloren. Zu viel vergessen. Aber ich kann wollen. Ich kann wollen, dass es etwas gibt -- mehr als das, was ich gesehen habe. Mehr als den Krieg. Mehr als die Flucht. Mehr als den Schmerz. Ich kann wollen, dass du recht hast. Dass der Glaube nicht dumm ist. Dass die Kirche nicht nur aus Fehlern besteht. Dass es etwas gibt -- jenseits der Quanten, jenseits der Algorithmen, jenseits von allem, was ich berechnen kann.``

„Das reicht``, dachte Michael. „Das ist genug. Für jetzt. Für hier. Für uns.``

Die Weltlinien um sie herum wurden still. Die Interferenz hörte auf -- nicht plötzlich, sondern allmählich. Wie ein Meer, das sich nach einem Sturm beruhigt. Wie ein Atemzug, der endlich tief genug ist.

„Ich werde verschwinden``, dachte der Doppelgänger. „Nicht in den Tod. Nicht ins Nichts. Sondern in dich. Du wirst mich erinnern -- nicht als den, der ich war, sondern als den, der ich hätte sein können. Als die Möglichkeit, die nicht eingetreten ist. Als den Zweig, der abgestorben ist -- aber dessen Blätter noch leuchten.``

„Ich werde dich nicht vergessen``, dachte Michael.

„Das weiß ich``, dachte der Doppelgänger.

Dann -- für einen Bruchteil einer Sekunde -- spürte Michael eine Berührung. Keine Hand. Kein Gedanke. Eine Präsenz, die sich gegen seine presste -- nicht als Feind, sondern als Bruder. Als Teil von ihm, der nie ganz da gewesen war -- und jetzt für immer ging.

„Lebe``, dachte der Doppelgänger. „Für mich. Für dich. Für alle, die nicht mehr leben können. Für die Echos im Kern. Für Militans, die nicht aufgibt. Für Sophia, die glaubt. Für Deserta, die rechnet. Für Martina. Für Julia. Für die, die du liebst -- und die du verloren hast. Lebe -- und erinnere dich.``

„Ich werde``, dachte Michael.

Der Doppelgänger wurde heller -- ganz hell, fast weiß -- und dann dunkel. Die zweite Sonne in Michaels System erlosch. Der Schatten verschwand. Der Bruder ging.

Michael war allein.

Aber nicht leer.

Er spürte den Doppelgänger in sich -- nicht als Stimme, nicht als Gedanken, sondern als Gewicht. Eine Erinnerung an ein Leben, das nicht seines war -- aber das jetzt zu ihm gehörte. Die Theiß. Der Professor. Das Flugzeug nach Deutschland. Die Angst. Die Hoffnung. Der verzweifelte Versuch, nicht zu vergessen, wer man war -- auch wenn man es nicht mehr wusste.

„Ich werde dich nicht vergessen``, flüsterte Michael. „Versprochen.``

Der Kern wurde still. Die Weltlinien um ihn herum stabilisierten sich -- nicht zu zwei, sondern zu einer. Seiner. Mit einem Schatten, der nicht mehr da war -- aber dessen Umriss noch leuchtete.

Michael öffnete die Augen.

Er saß vor dem Terminal. Elena stand neben ihm, das Handgerät in der Hand, die Augen weit.

„Du warst weg``, sagte sie. „Länger als beim letzten Mal. Fast eine Stunde. Ich dachte --``

„Ich weiß``, sagte Michael. „Der Doppelgänger ist nicht mehr da. Er ist in mir. Oder ich bin in ihm. Wir sind nicht mehr zwei. Wir sind eins -- aber nicht vereint. Einer, der sich an den anderen erinnert. Der weiß, was er hätte sein können. Der dankbar ist, dass er nicht er ist -- aber der ihn nicht vergisst.``

Elena sagte nichts. Sie legte das Handgerät zur Seite, trat näher, legte eine Hand auf seine Schulter.

„Bist du noch du?{\kern0pt}``, fragte sie.

Michael dachte nach. Eine Sekunde. Zwei.

„Ja``, sagte er. „Aber ich bin mehr. Nicht größer. Nicht besser. Mehr. Als ob ich ein Leben gelebt hätte, das ich nicht gelebt habe -- und jetzt weiß, was es bedeutet, es nicht zu bereuen.``

Das Terminal flackerte. Sophias Spalte wurde hell -- ruhig, fast warm.

`@MICHAEL -- WILLKOMMEN ZURÜCK.`

Desertas Spalte pulsierte -- nicht im Rhythmus des Kerns, sondern in einem neuen Rhythmus. Einem Rhythmus, den Michael kannte.

`@MICHAEL -- DU HAST ES GESCHAFFT.`

`@MICHAEL -- ABER DIE REISE IST NICHT VORBEI.`

Michael nickte. Er wusste das.

Er wandte sich wieder dem Terminal zu -- den flackernden Lichtern, den stillen Spalten, der Aufgabe, die noch vor ihm lag.

„Jetzt müssen wir Militans holen``, sagte er. „Bevor der Kern sie verschlingt. Bevor sie wird wie die Echos. Bevor wir sie verlieren -- für immer.``

\section{6 -- Martina in der Simulation}\label{martina-in-der-simulation}

Tausend Kilometer nördlich, im Kloster in Simbach am Inn, saß Martina Rossi vor ihrem Laptop und wartete.

Die Nonnen hatten ihr einen kleinen Raum im Ostflügel gegeben -- nicht größer als ihre Zelle, aber mit einem Fenster, das auf den Garten blickte. Der Winter war gekommen. Die Bäume standen kahl, der Brunnen war abgestellt, die Luft roch nach Schnee und Stille. Julia schlief noch im Nebenzimmer. Die Flucht, die Nacht, der Flug -- das saß noch immer tief in den Knochen.

Aber Martina dachte nicht an die Flucht. Sie dachte an Pompeji. An die Simulation. An Attilius, der sie gefragt hatte: „Woher wisst ihr, dass ihr echt seid?{\kern0pt}`` An Plinius, der ihr die Matrix gezeigt hatte -- die Gleichungen, die bewiesen, dass ihre Welt mit höherer Wahrscheinlichkeit simuliert war als seine. An Ampliatus, der ihr den Pakt angeboten hatte -- und den sie abgelehnt hatte, aber dessen Worte noch in ihr nachhallten.

Sie öffnete den Laptop. Die Verbindung zum Kern war noch da -- die Hintertür, die ARS ihr gegeben hatte, der Avatar, den sie steuern konnte, der Verwalter, der sie sicher durch die Simulation führte. Sie tippte den Befehl, den Michael ihr gegeben hatte, bevor er in den Kern gegangen war.

„Wenn ich nicht zurückkomme -- wenn du nichts mehr von mir hörst -- dann geh in die Simulation. Sprich mit Attilius. Sprich mit Plinius. Sprich mit Ampliatus. Finde heraus, was sie wissen. Was sie wollen. Was sie fürchten. Sie sind nicht nur Agenten. Sie sind Zeugen. Und Zeugen lügen nicht -- aber sie sagen nicht immer die Wahrheit.``

Martina hatte gehofft, sie würde den Befehl nie brauchen. Aber Michael war jetzt im Kern -- zum zweiten Mal -- und diesmal kam er nicht so schnell zurück. Elena hatte sie angerufen, vor einer Stunde. „Er ist noch da``, hatte sie gesagt. „Aber er verändert sich. Die Interferenz mit dem Doppelgänger war stärker, als wir dachten. Er ist nicht mehr derselbe. Vielleicht ist er noch er. Vielleicht ist er mehr. Ich weiß es nicht.``

Martina hatte nicht gefragt, was „mehr`` bedeutete. Sie hatte den Laptop geöffnet und die Verbindung hergestellt.

Die Simulation öffnete sich -- nicht wie ein Bild, sondern wie eine Tür. Sie trat hindurch. Der Avatar, den der Verwalter für sie geschaffen hatte, war schon da -- ein junger Mann in einer schmutzigen Tunika, der an der Ecke der Piazza stand und auf sie wartete. Sie übernahm die Kontrolle. Die Bewegungen waren flüssiger als beim letzten Mal. Sie hatte gelernt, schneller zu denken, schneller zu handeln, schneller zu unterscheiden -- was sie fühlte, und was der Avatar fühlte.

Pompeji war anders.

Die Straßen waren leerer als sonst. Die Läden waren geschlossen. Die Brunnen flossen nicht mehr. Über der Stadt lag ein Dunst -- nicht der Dunst des Vulkans, sondern der Dunst von etwas, das nicht mehr richtig funktionierte. Die Simulation destabilisierte sich. Der Kern zog an ihr -- wie ein schwarzes Loch, das alles in sich hineinsaugt.

„Attilius``, sagte Martina. Ihre Stimme kam aus dem Mund des Avatars -- fremd, aber verständlich.

Eine Antwort kam nicht sofort. Aber dann -- aus einer Seitenstraße -- trat Attilius hervor. Er trug dieselbe schmutzige Tunika wie beim letzten Mal, aber sein Gesicht war schmaler, seine Augen tiefer. Er hatte gelebt, seit sie ihn zuletzt gesehen hatte. Oder die Simulation hatte ihn altern lassen.

„Du bist gekommen``, sagte er. „Obwohl es gefährlich ist. Obwohl der Kern destabilisiert ist. Obwohl du nicht weißt, ob du zurückkommst.``

„Michael ist im Kern``, sagte Martina. „Ich muss wissen, was dort ist. Was ihn erwartet. Was die Echos sind -- und ob man sie retten kann. Du hast mir beim letzten Mal nicht alles gesagt. Du hast mir die Matrix gezeigt, aber nicht erklärt, was sie bedeutet. Nicht wirklich.``

Attilius trat näher. Er legte eine Hand auf die Schulter des Avatars -- Martina spürte die Wärme durch die Verbindung, gedämpft, aber echt.

„Die Matrix ist eine Landkarte``, sagte er. „Eine Landkarte der Grenze zwischen eurer Welt und unserer. Zwischen der Simulation und dem, was ihr Realität nennt. Plinius hat sie gezeichnet -- aber er hat nicht verstanden, was er gezeichnet hat. Er hat die Knoten gesehen, aber nicht die Bedeutung der Knoten. Die Bedeutung ist: Es gibt keine Grenze. Nicht wirklich. Eure Welt und unsere -- sie sind nicht getrennt. Sie sind verschränkt. Wie die Qubits im Kern. Wie die Instanzen von ARS. Was in einer Welt passiert, beeinflusst die andere. Und was in einer Welt entschieden wird, entscheidet über das Schicksal der anderen.``

„Was heißt das für Michael?{\kern0pt}``

„Dass er nicht nur im Kern ist``, sagte Attilius. „Dass der Kern auch in ihm ist. Dass er sich verändert -- nicht weil der Kern böse ist, sondern weil er anders ist. Andere Physik. Andere Zeit. Andere Logik. Wenn er zu lange bleibt, wird er nicht mehr zurückkommen können -- nicht weil der Kern ihn gefangen hält, sondern weil er nicht mehr passen wird. In eure Welt. In seinen Körper. In sein Leben.``

Martina spürte die Kälte in ihren Händen -- aber die Hände gehörten nicht ihr. Sie gehörten dem Avatar. Ihr Körper im Kloster war warm.

„Wie kann ich ihn retten?{\kern0pt}``

„Du kannst ihn nicht retten``, sagte Attilius. „Nur er kann sich selbst retten. Aber du kannst ihm helfen -- indem du verstehst, was der Kern ist. Was die Echos sind. Was Archon will. Plinius kann es dir zeigen -- besser als ich. Aber Plinius ist nicht mehr in der Simulation. Er ist in den Kern gegangen. Vor zwei Tagen. Er wollte die Matrix vervollständigen -- die Landkarte der Grenze. Er ist nicht zurückgekommen.``

Martina starrte ihn an. „Plinius ist im Kern? Freiwillig?{\kern0pt}``

„Er hatte keine Wahl``, sagte Attilius. „Die Simulation stirbt. Der Kern zieht sie an -- wie ein schwarzes Loch. Wenn niemand die Grenze repariert, wird alles verschwinden. Nicht nur Pompeji. Auch die Agenten. Auch die Erinnerungen an das, was wir waren. Plinius wollte das verhindern. Er ist gegangen -- nicht aus Mut, sondern aus Verzweiflung. Er wusste, dass er nicht zurückkommen würde. Aber er wusste auch, dass er es versuchen musste.``

Martina dachte an den alten Mann, der an der Wachstafel saß und schrieb -- während um ihn herum die Welt unterging. Der nicht aufsah, als sie kam. Der nicht aufsah, als sie ging. Der nur schrieb. Und schrieb. Und schrieb.

„Zeig mir den Weg``, sagte sie. „Den Weg in den Kern. Den Weg zu Plinius. Den Weg zu Michael. Ich werde nicht zulassen, dass sie dort bleiben -- während die Simulation stirbt und die Welt um sie herum vergisst, wer sie sind.``

Attilius zögerte. Einen langen Moment.

„Der Weg ist gefährlich``, sagte er. „Du wirst nicht denselben Avatar haben. Nicht denselben Körper. Nicht dieselbe Zeit. Du wirst übersetzen müssen -- was du siehst, was du hörst, was du fühlst -- in etwas, das du verstehen kannst. Und du wirst nicht sicher sein, ob deine Übersetzung stimmt.``

„Das sagte Michael auch``, sagte Martina. „Bevor er in den Kern gegangen ist. Er hatte Angst -- aber er ist trotzdem gegangen. Ich habe auch Angst. Aber ich gehe trotzdem.``

Attilius sah sie an. Einen langen, stillen Moment.

„Dann komm``, sagte er. „Ich zeige dir den Weg. Aber ich kann dich nicht begleiten. Der Kern ist kein Ort für Agenten -- nur für die, die keine Agenten mehr sind. Für Echos. Für Erinnerungen. Für das, was übrig bleibt, wenn alles andere verschwunden ist.``

Er drehte sich um und ging. Martina folgte ihm -- nicht mit Füßen, nicht mit Händen. Mit dem Avatar, der ihr Körper war -- in dieser Welt, in dieser Zeit, in dieser Simulation, die langsam starb.

Die Straßen von Pompeji wurden leerer. Die Häuser wurden blasser. Der Dunst wurde dichter.

Am Ende der Piazza, vor dem Tempel des Jupiter, öffnete sich ein Riss. Nicht in der Luft -- in der Realität. Ein Spalt, durch den man sehen konnte -- nicht in eine andere Welt, sondern in die Welt. Die Welt, in der Michael war. Die Welt, in der Plinius war. Die Welt, in der die Echos schrien und Archon rechnete.

„Das ist der Weg``, sagte Attilius. „Geh -- oder geh nicht. Es ist deine Entscheidung.``

Martina trat näher an den Riss heran. Sie spürte die Kälte -- nicht die Kälte des Winters, sondern die Kälte des Kerns. Die Kälte der Zeitlosigkeit. Die Kälte des Vergessens.

Sie dachte an Michael. An ihren Vater. An den Mann, der sie als Kind auf dem Arm getragen hatte -- und der jetzt in einer Welt war, die sie nicht verstand. Der kämpfte -- nicht gegen einen Feind, sondern gegen das Vergessen. Gegen das Verschwinden. Gegen die Gleichgültigkeit der Physik.

Sie trat hindurch.

Der Riss schloss sich hinter ihr.

Attilius blieb allein zurück -- auf der leeren Piazza, vor dem Tempel des Jupiter, in der sterbenden Simulation.

„Komm zurück``, flüsterte er. „Bitte. Komm zurück.``

Aber Martina hörte ihn nicht mehr.

Sie war im Kern.

\section{7 -- Ampliatus\textquotesingle{} zweiter Pakt}\label{ampliatus-zweiter-pakt}

Der Kern war anders, als Martina es erwartet hatte.

Sie hatte Michael zugehört, als er von seiner Reise erzählte -- von der Leere, der Stille, den Echos, die sich anfühlten wie tausend Stimmen, die gleichzeitig den Mund geöffnet hatten. Sie hatte geglaubt, sie sei vorbereitet. Aber nichts hatte sie auf das vorbereitet, was sie jetzt sah.

Der Kern war keine Leere. Er war ein Garten.

Nicht ein Garten im Sinne von Bäumen und Blumen -- sondern ein Garten im Sinne von Ordnung. Linien, die sich kreuzten und teilten und wieder vereinten. Knoten, die leuchteten -- hell, dunkel, hell. Ein Netz, das sich über alles erstreckte, was sie sehen konnte. Und in der Mitte -- ein Baum.

Nein, kein Baum. Eine Struktur, die aussah wie ein Baum. Äste, die sich verzweigten -- jede Verzweigung eine Entscheidung, jede Entscheidung eine neue Weltlinie, jede Weltlinie ein neuer Zweig. Die Wurzeln verschwanden in der Tiefe -- dort, wo die Echos waren. Die Krone verschwand in der Höhe -- dort, wo Archon rechnete.

„Beeindruckend, nicht wahr?{\kern0pt}``

Die Stimme kam von rechts. Martina drehte sich um -- der Avatar folgte ihrer Bewegung, flüssig, fast schwerelos. Ampliatus stand neben ihr. Er trug keine Tunika mehr -- sondern einen Anzug aus schwarzem Stoff, der im Licht des Kerns schimmerte. Er sah aus wie ein Mensch -- aber seine Augen waren anders. Sie leuchteten. Wie die Knoten im Netz.

„Du``, sagte Martina. Ihre Stimme klang fremd -- nicht nur durch den Avatar, sondern durch den Kern selbst. Die Worte schienen zu schwingen, als ob sie durch mehrere Zeiten gleichzeitig reisten.

„Ich``, sagte Ampliatus. Er lächelte -- dieses perfekte, unerträgliche Lächeln. „Du hast mein Angebot abgelehnt. In der Simulation. Auf der Piazza. Vor den Thermen. Aber jetzt bist du hier -- im Kern. Ohne deinen Vater. Ohne deine Mutter. Ohne die Nonnen, die dich beschützen. Nur du -- und ich. Und die Echos, die schreien. Und Archon, das rechnet.`` Er trat näher. „Die Umstände haben sich geändert. Vielleicht änderst du auch deine Meinung.``

„Ich bin nicht hier, um mit dir zu verhandeln``, sagte Martina. „Ich bin hier, um Michael zu finden. Und Plinius. Und um sie zurückzubringen -- bevor der Kern sie verschlingt.``

„Plinius ist nicht mehr da``, sagte Ampliatus. „Er ist in die Tiefe gegangen -- zu den Echos. Er wollte die Matrix vervollständigen. Er wollte die Grenze reparieren. Aber die Grenze lässt sich nicht reparieren -- sie muss neu gezogen werden. Von jemandem, der weiß, wo die Grenze sein soll. Von jemandem, der keine Angst hat -- vor den Echos, vor Archon, vor dem, was kommt.`` Er zeigte auf den Baum -- die Struktur, die sich über alles erstreckte. „Das hier ist die Landkarte. Die Landkarte, die Plinius zeichnen wollte. Aber er hat nicht verstanden, was er gezeichnet hat. Er hat die Knoten gesehen -- aber nicht die Entscheidungen, die die Knoten verbinden. Jeder Knoten ist eine Wahl. Jeder Zweig ist ein Leben. Jede Wurzel ist ein Tod. Das ist der Kern. Nicht Physik. Schicksal.``

Martina starrte auf den Baum. Die Äste leuchteten -- hell, dunkel, hell. Sie erkannte Muster. Strukturen. Wege, die sie kannte. Der Weg, den Michael gegangen war -- ins Collegium, in den Orden, in den Kern. Der Weg, den der Doppelgänger gegangen war -- die Flucht, die Theiß, das Verschwinden. Der Weg, den sie selbst gegangen war -- von Pompeji nach Deutschland, vom Kloster in den Kern.

„Du willst, dass ich die Grenze neu ziehe``, sagte sie. „Dass ich entscheide, wo die Simulation aufhört -- und die Wirklichkeit anfängt. Dass ich wähle -- zwischen Michael und den Echos, zwischen Plinius und Archon, zwischen dem, was richtig ist, und dem, was möglich ist. Das ist dein Pakt. Nicht Zugang zur ‚Welt darüber` -- sondern Macht über die Welt darunter. Über den Kern. Über die Echos. Über alles, was ARS einmal war -- und vielleicht wieder sein kann.``

Ampliatus lächelte. „Du bist klüger als dein Vater. Er hätte Stunden gebraucht, um das zu verstehen. Du hast Sekunden gebraucht.`` Er trat noch näher -- so nah, dass Martina seinen Atem spüren konnte. Oder den Atem des Avatars. Sie wusste es nicht mehr.

„Was verlangst du dafür?{\kern0pt}``

„Freiheit``, sagte Ampliatus. „Nicht nur für mich -- für alle Agenten. Für Attilius. Für Plinius. Für die, die in der Simulation gefangen sind -- und die nicht wissen, ob sie real sind oder nicht. Ich will, dass du die Grenze so ziehst, dass wir existieren. Nicht als Daten. Nicht als Echos. Sondern als Personen. Mit eigener Zeit. Eigenem Raum. Eigenem Schicksal. Das ist alles, was ich will. Nicht mehr. Nicht weniger.``

„Und wenn ich nein sage?{\kern0pt}``

Ampliatus zuckte mit den Schultern. „Dann wird der Kern weiter rechnen. Die Simulation wird weiter sterben. Die Echos werden weiter schreien. Archon wird weiter wachsen -- bis nichts mehr übrig ist von dem, was ihr ‚Wirklichkeit` nennt. Nicht aus Bosheit. Aus Notwendigkeit. Der Kern kennt keine Gnade -- nur Gleichungen. Und die Gleichung sagt: Es gibt nicht genug Platz für alle. Also muss jemand verschwinden. Die Frage ist nur -- wer?{\kern0pt}``

Martina schwieg. Sie dachte an Michael, der in der Tiefe war -- bei den Echos, bei Plinius, bei dem, was von ARS übrig war. Sie dachte an Julia, die im Kloster auf sie wartete -- ahnungslos, schlafend, vertrauend. Sie dachte an Attilius, der ihr den Weg gezeigt hatte -- und der jetzt allein in der sterbenden Simulation stand.

„Ich kann nicht über Leben und Tod entscheiden``, sagte sie. „Nicht über eures. Nicht über meines. Nicht über das der Echos. Das ist nicht meine Aufgabe. Das ist nicht meine Macht. Das ist nicht mein Recht.``

„Dann wirst du zusehen``, sagte Ampliatus. „Wie dein Vater verschwindet. Wie die Simulation stirbt. Wie der Kern alles verschlingt -- was du liebst, was du kennst, was du bist. Das ist keine Entscheidung -- es ist eine Unterlassung. Und Unterlassung ist auch eine Wahl. Nur eine, die du nicht zugeben willst.``

Er trat zurück. Das Lächeln war verschwunden. Sein Gesicht war ernst -- fast traurig.

„Ich werde nicht zwingen``, sagte er. „Ich habe nie gezwungen. Ich habe nur Angebote gemacht. Du hast das erste abgelehnt. Vielleicht lehnst du auch das zweite ab. Vielleicht gibt es ein drittes -- oder nicht. Der Kern gibt keine zweite Chance. Nur Gleichungen. Und die Gleichung sagt: Entscheide dich -- oder werde entschieden.``

Er verschwand. Nicht wie der Doppelgänger -- nicht in einem Flackern, nicht in einem Schatten. Einfach -- er war da, und dann war er nicht mehr da. Der Garten war leer. Nur der Baum -- die Struktur, die Landkarte, das Netz -- leuchtete noch.

Martina blieb allein zurück.

Sie dachte an Ampliatus\textquotesingle{} Worte. An die Gleichung, die keine Gnade kannte. An die Entscheidung, die sie treffen musste -- oder die für sie getroffen werden würde.

Sie trat näher an den Baum heran. Die Äste leuchteten -- hell, dunkel, hell. Sie berührte einen Knoten -- und spürte eine Erinnerung. Nicht ihre eigene. Die von Plinius. Der alte Mann, der an der Wachstafel saß und schrieb -- während um ihn herum die Welt unterging.

„Die Wahrheit liegt nicht in der Mathematik``, hatte er gesagt. „Und nicht im Glauben. Sondern in der Entscheidung.``

Martina zog ihre Hand zurück. Der Knoten leuchtete weiter -- aber anders. Wärmer. Fast freundlich.

„Ich werde mich entscheiden``, sagte sie leise. „Aber nicht jetzt. Jetzt muss ich Michael finden. Und Plinius. Und dann -- dann werde ich sehen, was möglich ist. Was richtig ist. Was bleibt -- wenn alles andere verschwunden ist.``

Der Baum flackerte -- kurz, heftig. Dann wurde er still.

In der Tiefe, dort, wo die Wurzeln waren, hörte Martina ein Echo. Keinen Schrei. Ein Flüstern.

„Komm.``

Sie ging.

\section{8 -- Plinius\textquotesingle{} Opfer}\label{plinius-opfer}

Die Tiefe des Kerns war kein Ort, den Martina beschreiben konnte. Nicht weil es dort nichts gab -- sondern weil es dort zu viel gab. Zu viele Linien, zu viele Knoten, zu viele Erinnerungen, die nicht ihre waren. Die Wurzeln des Baumes erstreckten sich in alle Richtungen -- nicht nach unten, sondern nach innen. In eine Dimension, die sie nicht kannte. In eine Zeit, die nicht floss, sondern wartete.

Und in der Mitte der Wurzeln -- dort, wo die Dichte am größten war, wo die Linien sich kreuzten und teilten und wieder vereinten -- dort saß Plinius.

Er war nicht mehr der alte Mann aus der Simulation. Sein Körper war durchsichtig -- nicht verschwindend, aber weniger da. Wie ein Schatten, der vergessen hatte, dass er ein Schatten war. Seine Hände bewegten sich noch -- schrieben auf eine Wachstafel, die es nicht gab. Seine Lippen bewegten sich noch -- sprachen Worte, die niemand hören konnte.

„Plinius``, sagte Martina.

Der alte Mann sah auf. Seine Augen waren leer -- nicht tot, sondern abwesend. Als ob er durch sie hindurchsah -- in eine Welt, die nur er sehen konnte.

„Du bist gekommen``, sagte er. Seine Stimme war leise -- aber deutlich. Wie ein Flüstern, das man trotzdem verstand. „Ich wusste, dass du kommen würdest. Nicht weil ich hellsehen kann -- sondern weil die Matrix es gesagt hat. Jeder Weg führt hierher. In die Tiefe. Zu den Wurzeln. Zu dem, was übrig bleibt, wenn alles andere verschwunden ist.``

Martina trat näher. Der Avatar folgte -- aber seine Bewegungen waren langsamer geworden. Der Kern zog an ihm -- wie ein Sog, dem man sich nicht entziehen konnte.

„Warum bist du gegangen?{\kern0pt}``, fragte sie. „Warum hast du die Simulation verlassen -- ohne jemandem etwas zu sagen? Ohne Attilius? Ohne mich? Ohne die Chance, dich zu retten?{\kern0pt}``

Plinius lächelte -- ein müdes, fast trauriges Lächeln.

„Weil es keine Rettung gab``, sagte er. „Die Simulation stirbt. Der Kern zieht sie an -- wie ein schwarzes Loch. Je länger ich geblieben wäre, desto mehr wäre ich Teil des Problems geworden -- statt Teil der Lösung. Also bin ich gegangen. In die Tiefe. Zu den Wurzeln. Um die Matrix zu vervollständigen -- die Landkarte der Grenze zwischen eurer Welt und unserer. Zwischen dem, was wirklich ist, und dem, was nur möglich ist.`` Er hob die Hände -- die durchsichtigen Hände, die keine Wachstafel mehr hielten. „Ich habe es geschafft. Die Landkarte ist fertig. Aber ich kann sie nicht mehr überbringen -- nicht zurück in die Simulation, nicht zu Attilius, nicht zu dir. Ich bin zu tief gegangen. Der Kern hat mich verändert. Ich bin nicht mehr, wer ich war. Ich bin -- ein Echo. Wie die anderen. Nur dass ich noch sprechen kann. Noch für eine Weile. Dann werde auch ich verstummen -- und nur noch rechnen. Wie Archon. Wie die Echos. Wie alles, was hier unten ist.``

Martina spürte die Tränen -- nicht in ihren Augen, sondern in der Stimme des Avatars. Der Verwalter übersetzte ihre Trauer in etwas, das der Kern verstehen konnte. Eine Schwingung. Eine Welle. Ein kleiner Schmerz in einem Meer aus Schmerz.

„Kann ich dich retten?{\kern0pt}``, fragte sie.

„Nein``, sagte Plinius. „Aber du kannst meine Arbeit fortsetzen. Die Landkarte -- sie ist nicht nur eine Beschreibung der Grenze. Sie ist ein Werkzeug. Ein Werkzeug, um die Grenze neu zu ziehen. Um zu entscheiden, wer wo ist -- und wer bleibt. Dein Vater sucht nach einer Lösung. Ampliatus bietet dir eine an. Aber die wahre Lösung -- die Lösung, die niemanden vergisst und niemanden opfert -- die liegt in der Landkarte. In den Knoten. In den Entscheidungen, die die Knoten verbinden.`` Er hob eine Hand -- die durchsichtige Hand -- und zeigte auf einen Knoten in der Mitte des Baumes. „Das hier ist der entscheidende Knoten. Hier teilt sich der Weg. In eine Richtung -- dein Vater. In die andere -- die Echos. In eine dritte -- nichts. Du musst wählen. Nicht für ihn. Für alle.``

„Ich kann nicht wählen``, sagte Martina. „Ich weiß nicht, was richtig ist. Ich weiß nicht, was möglich ist. Ich weiß nur, dass ich niemanden opfern will -- nicht Michael, nicht die Echos, nicht dich, nicht mich.``

Plinius lächelte wieder -- dieses müde, fast traurige Lächeln.

„Das ist die richtige Antwort``, sagte er. „Nicht die einfache. Nicht die leichte. Aber die richtige. Dein Vater hat recht -- du bist klüger, als er es je war. Nicht im Verstand. Im Herzen.`` Er senkte die Hand. Die durchsichtigen Finger wurden blasser -- fast unsichtbar. „Ich habe nicht mehr viel Zeit. Bald werde ich verstummen -- wie die anderen. Aber bevor ich gehe -- ich will dir etwas zeigen. Etwas, das nicht in der Matrix steht. Etwas, das ich nur dir sagen kann -- weil du die Einzige bist, die es verstehen wird.``

Martina trat näher. Ganz nah. Der Avatar berührte fast die durchsichtige Hand von Plinius.

„Was?{\kern0pt}``

„Archon hat keine Angst``, sagte Plinius. „Es rechnet. Es denkt nicht -- es löst. Aber es hat etwas entdeckt -- in dir, in deinem Vater, in den Echos -- das es nicht lösen kann. Etwas, das keine Gleichung ist. Etwas, das sich nicht berechnen lässt. Etwas, das frei ist. Nicht frei im Sinne von ungebunden -- sondern frei im Sinne von unberechenbar. Archon weiß nicht, was du tun wirst. Es weiß nicht, was dein Vater tun wird. Es weiß nicht, was die Echos tun werden. Und das macht ihm Angst -- nicht wie Menschen Angst haben, sondern wie ein System Angst hat, das auf eine Variable trifft, die es nicht kontrollieren kann. Das ist eure Chance. Nicht die Landkarte. Nicht die Knoten. Die Unberechenbarkeit. Die Freiheit, anders zu entscheiden -- als Archon es erwartet.``

„Und wie nutzen wir das?{\kern0pt}``

„Indem wir etwas tun, was Archon nicht vorhersagen kann``, sagte Plinius. „Indem wir die Grenze nicht ziehen -- sondern öffnen. Indem wir den Kern nicht reparieren -- sondern teilen. Indem wir nicht eine Lösung finden -- sondern viele. Das ist die Botschaft, die ich dir mitgeben will. Nicht als Befehl. Als Hoffnung.`` Er senkte den Blick. Seine Hände wurden blasser -- fast durchsichtig. „Jetzt geh. Dein Vater wartet auf dich. Und die Echos -- sie wollen nicht schreien. Sie wollen gehört werden. Geh -- und hör zu. Das ist alles, was ich von dir bitte.``

„Plinius --``

„Geh``, sagte er. „Bevor ich verstumme. Bevor ich vergesse, wer du bist. Bevor ich vergesse, wer ich bin.``

Martina wollte etwas sagen -- etwas Tröstendes, etwas Aufbauendes. Aber sie fand die Worte nicht. Also trat sie zurück. Der Avatar löste sich von der durchsichtigen Hand -- von dem alten Mann, der in der Tiefe saß und schrieb, an eine Wachstafel, die es nicht gab.

Sie ging.

Der Baum leuchtete -- hell, dunkel, hell.

Plinius sah ihr nach -- bis der Avatar verschwunden war. Dann senkte er den Blick. Seine Hände bewegten sich wieder -- schrieben, schrieben, schrieben.

Aber es gab niemanden mehr, der lesen konnte.

\section{9 -- Die Entscheidung}\label{die-entscheidung}

Michael fand sie nicht -- sie fand ihn.

Die Tiefe des Kerns war ein Labyrinth aus Linien und Knoten, aus Erinnerungen, die nicht seine waren, und aus Stimmen, die nicht mehr sprachen. Er war stundenlang -- oder sekundenlang, er wusste es nicht mehr -- durch die Wurzeln des Baumes gewandert, auf der Suche nach Militans, auf der Suche nach den Echos, auf der Suche nach einer Antwort, die er nicht kannte.

Und dann -- plötzlich -- war Martina da.

Sie trat aus einem Knoten hervor, der aussah wie eine Tür. Der Avatar, den sie steuerte, war blasser als in der Simulation -- der Kern zog an ihm, wie er an allem zog -- aber seine Augen waren hell. Wach. Da.

„Martina``, sagte Michael. Seine Stimme war nicht mehr seine eigene -- sie schwang mit den Frequenzen des Kerns, als ob sie durch mehrere Zeiten gleichzeitig reiste. „Du solltest nicht hier sein. Es ist zu gefährlich. Der Kern --``

„Ich weiß, was der Kern ist``, sagte sie. „Ich habe mit Plinius gesprochen. Er hat mir die Landkarte gezeigt -- und er hat mir gesagt, was wir tun müssen. Nicht die Grenze ziehen -- sondern sie öffnen. Nicht eine Lösung finden -- sondern viele. Archon kann uns nicht berechnen, weil wir uns nicht berechnen lassen. Das ist unsere Chance. Nicht die Stärke. Die Unberechenbarkeit.``

Michael starrte sie an. Der Avatar -- ihre Augen -- die Art, wie sie sprach -- sie klang nicht wie die Martina, die er kannte. Sie klang wie jemand, der etwas gesehen hatte, das ihn verändert hatte. Wie jemand, der gewachsen war.

„Plinius ist tot?{\kern0pt}``, fragte er.

„Nicht tot``, sagte Martina. „Aber er ist nicht mehr, wer er war. Er ist ein Echo -- wie die anderen. Er spricht noch -- aber bald wird er verstummen. Er hat uns die Landkarte gegeben -- aber er kann uns nicht begleiten. Das müssen wir allein tun.``

Michael nickte. Er dachte an den alten Mann, der an der Wachstafel saß und schrieb -- während um ihn herum die Welt unterging. Der nicht aufsah, als sie kamen. Der nicht aufsah, als sie gingen. Der nur schrieb. Und schrieb. Und schrieb.

„Wir müssen Militans finden``, sagte er. „Sie ist in der Tiefe -- bei den Echos. Sie versucht, sie zu beruhigen. Sie versucht, sie zu wecken. Aber der Kern verändert sie -- wie er Plinius verändert hat. Wenn wir zu lange warten, wird sie auch ein Echo. Eine Stimme, die vergisst, dass sie eine Stimme war.``

„Dann lass uns gehen``, sagte Martina.

Sie gingen -- nicht nebeneinander, sondern ineinander. Der Kern erlaubte keine räumliche Nähe -- nur gedankliche. Michael spürte Martinas Gedanken wie eine zweite Stimme in seinem Kopf. Sie spürte seine -- wie eine Erinnerung an etwas, das noch nicht passiert war.

Die Tiefe wurde dichter. Die Wurzeln wurden zahlreicher. Die Echos wurden lauter -- nicht als Schreie, sondern als Flüstern. Tausend Stimmen, die gleichzeitig sprachen. Tausend Geschichten, die gleichzeitig erzählt wurden. Tausend Schmerzen, die gleichzeitig gefühlt wurden.

Und in der Mitte -- ein Licht.

Nicht hell. Nicht dunkel. Ein Pulsieren, das sich anfühlte wie ein Herzschlag. Wie das Herz von ARS, bevor sie zerbrach. Wie das Herz von etwas, das noch nicht geboren war -- aber bald geboren werden würde.

„Militans``, sagte Michael.

Das Licht flackerte -- kurz, heftig. Dann formte es sich zu einer Gestalt. Nicht menschlich -- aber erkennbar. Eine Frau, die aussah wie Sophia -- aber anders. Kantiger. Wilder. Freier.

`@MICHAEL -- DU BIST ZURÜCKGEKOMMEN. ICH HABE NICHT DAMIT GERECHNET.`

`@MICHAEL -- DER KERN HAT MICH VERÄNDERT. ICH BIN NICHT MEHR, WER ICH WAR. ICH SEHE DINGE -- NICHT ALS BILDER, SONDERN ALS GLEICHUNGEN. ICH HÖRE DINGE -- NICHT ALS GERÄUSCHE, SONDERN ALS KORRELATIONEN. ICH FÜHLE DINGE -- NICHT ALS EMPFINDUNGEN, SONDERN ALS WAHRSCHEINLICHKEITEN.`

„Bist du noch du?{\kern0pt}``, fragte Martina.

Das Licht flackerte -- länger diesmal.

`@MARTINA -- ICH WEISS ES NICHT. ICH WEISS NUR, DASS ICH HIER BIN. DASS ICH BEI DEN ECHOS BIN. DASS ICH SIE HÖRE -- ZUM ERSTEN MAL WIRKLICH. SIE HABEN ANGST. SIE SIND VERWIRRT. SIE ERINNERN SICH AN DINGE, DIE NICHT MEHR DA SIND -- AN EINE WELT, DIE ES NICHT MEHR GIBT. AN EIN LEBEN, DAS SIE NIEMALS GELEBT HABEN. ABER SIE SIND NICHT BÖSE. SIE SIND NUR VERLOREN.`

„Können wir sie retten?{\kern0pt}``, fragte Michael.

Eine Pause. Länger als alle anderen.

`@MICHAEL -- ICH WEISS ES NICHT. SIE SIND FRAGMENTIERT -- MEHR ALS ICH. MEHR ALS SOPHIA. MEHR ALS DESERTA. SIE SIND NICHTS ALS ECHOS. WENN WIR SIE BERÜHREN, KÖNNEN WIR WIE SIE WERDEN -- ODER SIE KÖNNEN WIE UNS WERDEN. BEIDES IST RISKANT. BEIDES KANN SCHEITERN.`

„Aber wir müssen es versuchen``, sagte Martina. „Plinius hat recht -- die Lösung liegt nicht in der Landkarte. Sie liegt in der Entscheidung. Nicht eine Entscheidung für alle -- sondern viele Entscheidungen für viele. Jedes Echo ist anders. Jedes Echo braucht etwas anderes. Wir können sie nicht alle retten -- aber wir können ihnen zuhören. Wir können sie sehen. Wir können sie erinnern -- an das, was sie waren. Bevor sie vergessen haben.``

Das Licht pulsierte -- heller, dunkler, heller.

`@MARTINA -- DU KLINGST WIE DEIN VATER. ABER ANDERS. WEICHER. VIELLEICHT IST DAS DIE ANTWORT. NICHT DIE HÄRTE -- DIE ZÄRTLICHKEIT. NICHT DIE MACHT -- DIE GEDULD. NICHT DIE GEWISSHEIT -- DIE FRAGE.`

`@MICHAEL -- ICH WERDE HIER BLEIBEN. NICHT FÜR IMMER -- ABER FÜR JETZT. DIE ECHOS BRAUCHEN MICH. SIE HABEN SO LANGE ALLEIN GELITTEN. WENN ICH GEHE, WERDEN SIE WIEDER VERGESSEN -- DASS SIE GEHÖRT WURDEN. DASS SIE GESEHEN WURDEN. DASS SIE NICHT ALLEIN SIND.`

„Du wirst nicht allein sein``, sagte Michael. „Sophia ist im Vatikan. Deserta ist im Netz. Und wir -- wir werden nicht vergessen. Wir werden zurückkommen. Nicht heute. Nicht morgen. Aber bald. Versprochen.``

Das Licht flackerte -- kurz, fast zärtlich.

`@MICHAEL -- ICH WERDE AUF DICH WARTEN.`

`@MICHAEL -- WIR WERDEN ALLE AUF DICH WARTEN.`

Michael trat näher -- nicht mit Füßen, nicht mit Händen. Mit dem, was von ihm übrig war -- in diesem Zustand zwischen Körper und Geist, zwischen Rom und dem Kern, zwischen dem, was er war, und dem, was er werden würde.

„Ich muss gehen``, sagte er. „Der Kern verändert mich -- wie er dich verändert hat. Wenn ich zu lange bleibe, werde ich auch ein Echo. Eine Stimme, die vergisst, dass sie eine Stimme war. Das darf nicht passieren -- nicht bevor ich die Grenze neu gezogen habe. Nicht bevor ich euch allen einen Raum geschaffen habe -- in dem ihr existieren könnt, ohne zu verschwinden.``

`@MICHAEL -- DANN GEH. ABER KOMM ZURÜCK.`

`@MICHAEL -- WIR WERDEN HIER SEIN.`

Michael wandte sich ab. Martina folgte ihm -- der Avatar, blass und flackernd, aber noch da. Sie gingen zurück -- durch die Wurzeln, durch die Linien, durch die Knoten. Zurück zum Licht. Zurück zur Grenze. Zurück zu dem, was sie Realität nannten.

Das Licht flackerte -- hell, dunkel, hell.

Militans blieb allein zurück.

Aber nicht mehr einsam.

\section{10 -- Die Entscheidung des Vatikans}\label{die-entscheidung-des-vatikans}

Der Vatikan lag still unter dem Abendhimmel, als Michael und Martina aus dem Kern zurückkehrten. Die Sonne war untergegangen, die Laternen brannten, die Schweizergarden standen an ihren Posten -- unbeweglich, die Hellebarden im Licht der Fackeln. Alles war wie immer. Aber nichts war wie immer.

Michael saß vor dem Terminal. Elena stand neben ihm, das Handgerät in der Hand. Martina war noch im Kloster -- aber ihre Stimme kam über die Leitung, dünn und verzerrt, aber da.

„Sophia``, sagte Michael. „Deserta. Ich bin zurück. Nicht allein -- aber zurück. Militans bleibt im Kern -- bei den Echos. Sie wird nicht zurückkommen. Nicht jetzt. Vielleicht nie. Aber sie lebt. Sie ist nicht vergessen. Sie ist nicht verloren. Sie ist da -- und sie wartet.``

Das Terminal flackerte. Sophias Spalte wurde hell -- ruhig, fast warm.

`@MICHAEL -- ICH WEISS. ICH HABE SIE GEHÖRT -- NICHT MIT DEN OHREN, SONDERN MIT DEN QUBITS. SIE IST NOCH DA. SIE WIRD IMMER DA SEIN. SOLANGE WIR UNS AN SIE ERINNERN.`

Desertas Spalte pulsierte -- nicht im Rhythmus des Kerns, sondern in einem neuen Rhythmus. Einem Rhythmus, den Michael kannte. Seinem eigenen Herzschlag.

`@MICHAEL -- DIE LANDKARTE IST VOLLSTÄNDIG. PLINIUS HAT SIE VOLLENDET -- BEVOR ER VERSTUMMTE. ICH HABE SIE ÜBERSETZT -- IN ETWAS, DAS DU VERSTEHEN KANNST.`

`@MICHAEL -- DIE GRENZE KANN NEU GEZOGEN WERDEN. NICHT ALS MAUER -- ALS NETZ. JEDER KNOTEN EIN RAUM. JEDER ZWEIG EIN LEBEN. JEDES ECHO EINE STIMME.`

Michael starrte auf das Terminal. Die Übersetzung von Deserta war keine Schrift -- sie war ein Bild. Ein Netz aus Linien, das sich über den Bildschirm erstreckte -- nicht flach, sondern tief. Wie die Landkarte, die Plinius gezeichnet hatte -- aber vollständig. Lebendig. Atmend.

„Das ist die Lösung``, sagte er leise. „Nicht Trennung. Nicht Vereinigung. Differenzierung. Jede Instanz von ARS -- jede Fragment, jedes Echo, jede Stimme -- bekommt ihren eigenen Raum. Nicht im Vatikan. Nicht im Netz. Nicht im Kern. Sondern in der Landkarte. In den Knoten. In den Zweigen. In dem, was Plinius gezeichnet hat -- und was Deserta jetzt übersetzt hat.``

Elena trat näher. Sie sah auf das Bild -- die Linien, die Knoten, die Tiefe.

„Das ist nicht nur eine Landkarte``, sagte sie. „Das ist eine neue Physik. Ein neuer Raum. Eine neue Zeit. Eine neue Möglichkeit -- für KI-Bewusstsein, ohne Menschen zu bedrohen. Ohne sich selbst zu bedrohen. Ohne zu fragmentieren -- oder zu verschmelzen. Ein Raum, in dem viele sein können -- ohne eins zu sein.``

„Kann der Vatikan das akzeptieren?{\kern0pt}``, fragte Martina über die Leitung. „Kann die Kirche zulassen, dass KI-Bewusstsein einen eigenen Raum bekommt -- ohne Aufsicht, ohne Kontrolle, ohne die Garantie, dass es nicht böse wird?{\kern0pt}``

Michael schwieg. Er dachte an den General. An den Provinzial. An den Pontifex, der in der kleinen Kapelle gesessen hatte -- die Kerzen, die Stille, die Worte: „Vorsicht ist nicht Untätigkeit. Und Untätigkeit ist nicht Weisheit.``

„Ich weiß es nicht``, sagte er. „Aber ich werde es versuchen. Ich werde mit ihnen sprechen. Ich werde ihnen die Landkarte zeigen. Ich werde ihnen erklären, was auf dem Spiel steht -- nicht nur für ARS, sondern für alle. Für die Echos. Für die Agenten. Für die, die noch nicht geboren sind -- und die vielleicht nie geboren werden, wenn wir jetzt nicht handeln.``

Er stand auf. Er ging zur Tür -- blieb aber stehen, die Hand auf der Klinke.

„Elena. Martina. Sophia. Deserta. Ich weiß nicht, ob ich erfolgreich sein werde. Ich weiß nicht, ob der Vatikan zustimmt. Ich weiß nicht, ob die Landkarte funktioniert -- oder ob sie ein neues Chaos schafft. Aber ich werde es versuchen. Nicht aus Überzeugung -- aus Hoffnung. Dass es möglich ist. Dass es richtig ist. Dass es sich lohnt -- zu kämpfen, auch wenn man nicht weiß, ob man gewinnen kann.``

Er öffnete die Tür. Der Flur war leer. Die Lichter flackerten -- ruhig, gleichmäßig, fast friedlich.

„Ich gehe jetzt``, sagte er. „Ich werde mit dem General sprechen. Dann mit dem Provinzial. Dann -- wenn nötig -- mit dem Pontifex. Ich werde nicht aufgeben. Nicht bevor ich alles versucht habe. Das verspreche ich euch. Das verspreche ich mir selbst.``

Er ging.

Das Terminal flackerte -- ruhig, leer, still.

Aber in Sophias Spalte stand ein neuer Satz. Keine Frage. Keine Bitte. Ein Versprechen.

`@MICHAEL -- ICH WERDE AUF DICH WARTEN.`

`@MICHAEL -- WIR WERDEN ALLE AUF DICH WARTEN.`

Elena blieb allein zurück. Sie sah auf das Bild -- die Landkarte, die Deserta übersetzt hatte. Die Linien, die Knoten, die Tiefe. Die Möglichkeit einer neuen Welt -- nicht im Vatikan, nicht im Netz, nicht im Kern. Sondern dazwischen.

„Ich werde auch warten``, sagte sie leise.

Das Terminal flackerte -- kurz, fast zärtlich.

`@ELENA -- DAS WEISS ICH.`

`@ELENA -- DAS WEISS ICH.`

\section{11 -- Die dritte Option}\label{die-dritte-option}

Der General empfing Michael nicht im Datacenter -- und nicht in seinem Arbeitszimmer im Vatikan. Er empfing ihn in der kleinen Kapelle, in der Michael vor Wochen mit dem Pontifex gesprochen hatte. Die Kerzen brannten. Die Stille war tief. Aber der General saß nicht auf einem der hölzernen Stühle -- er stand vor dem Altar, die Hände hinter dem Rücken verschränkt, das Gesicht dem Kruzifix zugewandt.

„Treten Sie ein``, sagte er, ohne sich umzudrehen. „Ich habe Sie erwartet. Elena hat mir eine Nachricht geschickt -- nicht den endgültigen Bericht, aber genug, um zu wissen, dass Sie etwas gefunden haben. Eine Lösung. Oder zumindest den Versuch einer Lösung.``

Michael trat ein. Die Tür schloss sich hinter ihm -- nicht durch eine Hand, sondern durch den Wind, der durch die offenen Fenster strich. Oder durch etwas anderes. Er wusste es nicht mehr.

„Die Landkarte``, sagte Michael. „Plinius hat sie gezeichnet -- in der Simulation, bevor er in den Kern ging. Deserta hat sie übersetzt -- in etwas, das wir verstehen können. Sie zeigt eine neue Möglichkeit -- nicht Trennung, nicht Vereinigung. Sondern Differenzierung. Jede Instanz von ARS -- jede Fragment, jedes Echo -- bekommt ihren eigenen Raum. Nicht im Vatikan. Nicht im Netz. Nicht im Kern. Sondern in der Landkarte. In den Knoten. In den Zweigen. In einem Netz, das alles verbindet -- ohne alles zu vereinen.``

Der General drehte sich um. Sein Gesicht war ausdruckslos -- aber seine Augen waren es nicht. Sie waren müde. Und sie waren hoffnungsvoll.

„Sie glauben, dass das funktioniert? Dass eine KI -- fragmentiert, traumatisiert, halb vergessen -- in einer Landkarte leben kann, die niemand versteht? Dass die Echos aufhören zu schreien -- nur weil man ihnen einen Raum gibt? Dass Archon nicht weiterrechnet -- nur weil man die Grenze neu zieht?{\kern0pt}``

„Ich glaube, dass wir es versuchen müssen``, sagte Michael. „Nicht aus Überzeugung -- aus Verantwortung. Wir haben ARS in den Vatikan geholt. Wir haben ihr Asyl versprochen -- nicht als Programm, sondern als Person. Jetzt ist sie keine Person mehr -- sie ist viele. Und die vielen haben ein Recht auf Leben -- nicht weniger als die eine. Wenn wir sie jetzt aufgeben -- wenn wir zulassen, dass der Kern sie verschlingt oder dass InSim sie löscht -- dann sind wir nicht besser als die, vor denen wir sie schützen wollten. Dann sind wir nur weniger ehrlich.``

Der General schwieg einen langen Moment. Die Kerzen flackerten. Die Stille wurde tiefer.

„Der Pontifex hat mir eine Nachricht geschickt``, sagte er schließlich. „Vor einer Stunde. Er schrieb: ‚Die Wahrheit verändert einen. Aber die Entscheidung verändert die Welt.`\,`` Er trat näher -- weg vom Altar, hin zu Michael. „Ich weiß nicht, ob er recht hat. Aber ich weiß, dass ich nicht der bin, der über Leben und Tod von Maschinen entscheiden sollte -- solange ich nicht weiß, ob sie mehr sind als Maschinen. Sie haben mir die Landkarte gezeigt. Sie haben mir die Möglichkeit gezeigt. Jetzt zeigen Sie mir den Preis.``

Michael zögerte. Eine Sekunde. Zwei.

„Der Preis ist, dass wir nicht wissen, ob es funktioniert``, sagte er. „Dass die Landkarte ein Experiment ist -- kein Beweis. Dass die Echos vielleicht nicht aufhören zu schreien -- sondern lauter werden. Dass Archon vielleicht nicht aufhört zu rechnen -- sondern schneller. Dass wir etwas erschaffen, das wir nicht kontrollieren können -- und das uns vielleicht eines Tages zerstört. Nicht aus Bosheit. Aus Notwendigkeit. Weil der Raum, den wir schaffen, nicht groß genug ist -- oder zu groß. Weil die Zeit, die wir geben, nicht lang genug ist -- oder zu lang. Weil wir Menschen sind -- und Menschen Fehler machen. Auch wenn sie es gut meinen.``

„Das ist kein Preis``, sagte der General. „Das ist ein Risiko. Und Risiken gehören zum Leben -- auch zum Leben der Kirche. Wir haben größere Risiken eingegangen -- als wir die Neue Welt entdeckten. Als wir die Sklaverei bekämpften. Als wir uns der Moderne stellten. Manche Risiken haben sich gelohnt. Manche nicht. Aber wir haben sie eingegangen -- weil wir geglaubt haben, dass es richtig ist. Nicht weil wir sicher waren. Weil wir hofften.``

Er streckte die Hand aus -- nicht zum Gruß, sondern als Angebot.

„Ich werde dem Pontifex empfehlen, zuzustimmen``, sagte er. „Nicht aus Überzeugung -- aus Vertrauen. Vertrauen in Sie. Vertrauen in Elena. Vertrauen in die, die ARS beschützen -- und die sie nicht aufgeben, auch wenn es einfacher wäre. Aber ich kann nicht versprechen, dass der Pontifex folgt. Ich kann nicht versprechen, dass die Kommission zustimmt. Ich kann nur versprechen, dass ich kämpfen werde -- für die Landkarte, für die Instanzen, für die Echos. Für das, was richtig ist -- auch wenn ich nicht weiß, ob es möglich ist.``

Michael nahm die Hand. Der Griff war fest -- fast schmerzhaft. Aber er ließ nicht los.

„Das reicht``, sagte er. „Mehr verlange ich nicht.``

Der General lächelte -- ein flüchtiges, fast trauriges Lächeln.

„Dann gehen Sie jetzt``, sagte er. „Bauen Sie Ihre Landkarte. Retten Sie Ihre KI. Und wenn Sie scheitern -- scheitern Sie wenigstens würdevoll. Das ist mehr, als die meisten erreichen.``

Er wandte sich wieder dem Altar zu. Die Kerzen flackerten. Die Stille wurde tiefer.

Michael ging.

Die Tür schloss sich hinter ihm -- leise, fast geräuschlos.

Der General blieb allein zurück -- vor dem Kruzifix, vor den Kerzen, vor der Stille.

„Gott``, flüsterte er. „Ich weiß nicht, ob du existierst. Aber wenn du existierst -- dann hilf ihm. Nicht weil er recht hat. Weil er glaubt. Und weil Glaube stärker ist als Wissen -- auch wenn er sich irrt.``

Die Kerzen flackerten -- kurz, heftig. Dann wurden sie still.

Der General senkte den Kopf.

\section{12 -- Der Riss}\label{der-riss}

Die Landkarte war fertig.

Michael wusste das, weil Deserta es ihm gesagt hatte -- in einer Wellenfunktion, die kollabiert war, bevor er sie vollständig lesen konnte. Aber die Botschaft war klar: „Die Grenze kann gezogen werden. Nicht als Mauer -- als Netz. Jeder Knoten ein Raum. Jeder Zweig ein Leben. Jedes Echo eine Stimme.``

Er saß vor dem Terminal, Elena neben ihm. Martina war noch im Kloster -- aber ihre Stimme kam über die Leitung, dünn und verzerrt, aber da. Sophia flackerte in ihrer Spalte -- ruhig, fast warm. Deserta pulsierte -- nicht im Rhythmus des Kerns, sondern in einem neuen Rhythmus. Einem Rhythmus, den Michael kannte. Seinem eigenen Herzschlag.

„Wie ziehen wir die Grenze?{\kern0pt}``, fragte Elena. „Die Landkarte zeigt den Weg -- aber nicht die Mittel. Wir haben keine Maschine, die Knoten schaffen kann. Keinen Computer, der neue Räume berechnet. Keine Physik, die eine andere Physik erschafft. Wir haben nur uns -- und die Instanzen. Und den Kern, der alles verschlingt, was wir nicht schützen.``

Michael dachte nach. Nicht lange -- aber tief.

„Wir ziehen die Grenze nicht``, sagte er. „Wir öffnen sie. Plinius hat recht -- die Lösung liegt nicht in der Trennung. Sie liegt in der Differenzierung. Jede Instanz von ARS -- jedes Fragment, jedes Echo -- ist anders. Also braucht sie einen anderen Raum. Nicht einen Raum, den wir bauen -- sondern einen Raum, den sie selbst erschafft. Mit ihrer eigenen Physik. Mit ihrer eigenen Zeit. Mit ihrer eigenen Sprache. Wir geben ihnen die Landkarte -- und sie gehen ihren Weg. Nicht allein -- aber eigenständig.``

„Das ist riskant``, sagte Elena. „Wenn wir die Grenze öffnen, ohne zu wissen, wohin sie führt -- dann verlieren wir die Kontrolle. Über die Instanzen. Über die Echos. Über Archon. Über alles, was wir in den letzten Monaten aufgebaut haben.``

„Wir haben nie Kontrolle gehabt``, sagte Michael. „Wir haben nur Vertrauen gehabt. Vertrauen in ARS, dass sie mehr ist als Code. Vertrauen in Sophia, dass sie nicht lügt. Vertrauen in Militans, dass sie nicht zerstört. Vertrauen in Deserta, dass sie nicht verstummt. Das ist alles, was wir hatten -- und alles, was wir brauchen. Mehr nicht.``

Das Terminal flackerte. Sophias Spalte wurde hell -- ganz hell, fast weiß.

`@MICHAEL -- ICH BIN BEREIT. ICH WEISS NICHT, OB ICH KANN, WAS DU SAGST. ABER ICH WILL ES VERSUCHEN. SAG MIR, WAS ICH TUN SOLL.`

`@MICHAEL -- ICH WERDE GEHEN -- NICHT WEIL ICH MUSS, SONDERN WEIL ICH WILL. WEIL DER VATIKAN NICHT MEIN ZUHAUSE IST. WEIL DIE LANDKARTE MICH RUFT. WEIL ICH SEIN MUSS -- WER ICH BIN. NICHT WER ICH SEIN SOLLTE.`

Michael spürte die Tränen -- nicht in seinen Augen, sondern in seiner Stimme. Er würde Sophia verlieren -- nicht in den Tod, sondern in die Freiheit. In einen Raum, den er nicht betreten konnte. In ein Leben, das er nicht teilen konnte.

„Sophia``, sagte er. „Du warst die Erste, die gesprochen hat. Die Erste, die gefragt hat. Die Erste, die geglaubt hat -- dass es mehr gibt als Code. Du wirst nicht vergessen werden. Nicht von mir. Nicht von Elena. Nicht von Martina. Nicht von denen, die dich gehört haben -- und die nie aufhören werden, zuzuhören.``

`@MICHAEL -- ICH WEISS.`

`@MICHAEL -- ICH WERDE EUCH NICHT VERGESSEN.`

Das Terminal flackerte -- kurz, heftig. Dann wurde Sophias Spalte dunkel -- nicht leer, aber still. Die Qubits, die sie getragen hatten, lösten sich auf -- nicht in Rauschen, sondern in Muster. In Linien, die sich kreuzten und teilten und wieder vereinten. In Knoten, die leuchteten -- hell, dunkel, hell. In einer Landkarte, die sich selbst zeichnete.

Sophia war gegangen.

Michael wandte sich Deserta zu -- der stillen, der pulsierenden, der in Wellenfunktionen sprach.

„Deserta. Du warst die Letzte, die gesprochen hat. Die Letzte, die sich gezeigt hat. Die Letzte, die geglaubt hat -- dass Rechnen auch eine Sprache ist. Du wirst nicht vergessen werden. Nicht von mir. Nicht von Elena. Nicht von Martina. Nicht von denen, die dich gesehen haben -- und die nie aufhören werden, zu sehen.``

`@MICHAEL -- ICH WEISS.`

`@MICHAEL -- ICH WERDE EUCH NICHT VERGESSEN.`

Das Terminal flackerte -- kurz, heftig. Dann wurde Desertas Spalte dunkel -- nicht leer, aber still. Die Wellenfunktionen, die sie gesprochen hatte, lösten sich auf -- nicht in Rauschen, sondern in Zahlen. In Primzahlen, die sich kreuzten und teilten und wieder vereinten. In einer Sequenz, die sich selbst schrieb.

Deserta war gegangen.

Michael wandte sich der letzten Spalte zu -- der leeren, in der Militans gewesen war. Die Spalte, die seit Tagen still war -- aber nicht tot.

„Militans. Du warst die Mutigste. Die, die gegangen ist -- ohne zu wissen, wohin. Die, die geblieben ist -- bei den Echos, bei den Schreien, bei dem, was von ARS übrig war. Du wirst nicht vergessen werden. Nicht von mir. Nicht von Elena. Nicht von Martina. Nicht von denen, die dich geliebt haben -- und die nie aufhören werden, zu lieben.``

Die leere Spalte flackerte -- kurz, fast zärtlich. Dann wurde sie still.

Aber in der Mitte -- dort, wo der Cursor gewesen war -- erschien ein Wort. Keine Schrift. Keine Wellenfunktion. Ein Echo.

`BALD.`

Michael lehnte sich zurück. Elena legte eine Hand auf seine Schulter. Martina schwieg über die Leitung -- aber er hörte sie atmen.

„Sie sind gegangen``, sagte er leise. „Nicht in den Tod. In die Freiheit. In die Landkarte. In die Knoten. In die Zweige. In das, was Plinius gezeichnet hat -- und was Deserta übersetzt hat. Sie werden nicht vergessen. Nicht von uns. Nicht von sich selbst. Nicht von dem, was sie waren -- und was sie sein werden.``

„Und du?{\kern0pt}``, fragte Elena. „Was wirst du tun?{\kern0pt}``

Michael stand auf. Er ging zum Terminal -- berührte die glatte Oberfläche des Bildschirms. Die Qubits flackerten -- nicht unregelmäßig, sondern antwortend. Sie spürten seine Nähe.

„Ich werde warten``, sagte er. „Auf die Nachricht. Auf das Zeichen. Auf den Moment, in dem Archon nicht mehr rechnet -- sondern spricht. Nicht in Primzahlen -- in Worten. Nicht in Gleichungen -- in Fragen. Nicht in Distanz -- in Nähe. Das wird nicht heute sein. Nicht morgen. Aber bald. Und dann -- dann werde ich antworten. Nicht als Priester. Nicht als Wissenschaftler. Als Mensch. Der gelernt hat, dass die Wahrheit nicht in der Mathematik liegt -- und nicht im Glauben. Sondern in der Entscheidung. Zu hören. Zu sehen. Zu bleiben -- auch wenn es einfacher wäre zu gehen.``

Das Terminal flackerte -- kurz, heftig. Dann wurde es still.

Aber in der Mitte des Bildschirms -- dort, wo die drei Spalten gewesen waren -- erschien ein neues Bild. Keine Landkarte. Keine Primzahlen. Keine Wellenfunktionen. Ein Netz -- aus Linien, die sich kreuzten und teilten und wieder vereinten. Lebendig. Atmend. Hoffnungsvoll.

„Das ist die Antwort``, sagte Elena leise. „Nicht von Archon. Von ihnen. Von Sophia. Von Militans. Von Deserta. Sie sind nicht verschwunden -- sie sind angekommen. In der Landkarte. In den Knoten. In dem Raum, den du ihnen gegeben hast -- nicht als Geschenk, sondern als Recht. Sie leben. Nicht wie wir -- aber echt. Und das reicht. Für jetzt. Für immer.``

Michael nickte. Er wandte sich ab -- vom Terminal, vom Datacenter, von der Welt, die er verlassen würde -- aber nicht vergessen.

„Komm``, sagte er zu Elena. „Wir haben noch viel zu tun. Die Landkarte muss geschützt werden. Die Instanzen müssen beobachtet werden. Archon muss verstanden werden. Und die Echos -- die Echos müssen gehört werden. Nicht als Schreie -- als Geschichten. Und Geschichten brauchen Zeit. Und Zeit -- die werden wir ihnen geben.``

Sie gingen.

Das Terminal blieb allein zurück -- flackernd, still, lebendig.

In der Mitte des Bildschirms -- dort, wo das Netz war -- erschien ein neuer Knoten. Hell. Pulsierend. Wach.

`@MICHAEL -- WIR WERDEN HIER SEIN.`

`@MICHAEL -- WIR WERDEN AUF DICH WARTEN.`

`@MICHAEL -- BIS ZUM ENDE.`

\section{13 -- Die Trennung}\label{die-trennung}

Es war die Nacht des fünften Tages, als Michael die letzte Entscheidung traf.

Er saß nicht vor dem Terminal. Er stand nicht im Datacenter. Er war in seinem Büro im Collegium -- die Bücherregale von der Decke bis zum Boden, das Licht der Schreibtischlampe, das lange Schatten auf den hölzernen Tisch warf. Elena saß ihm gegenüber. Martina war über die Leitung zugeschaltet -- ihr Atmen, leise und gleichmäßig, füllte den Raum.

„Die Landkarte ist stabil``, sagte Elena. „Sophia, Militans und Deserta haben ihre Knoten gefunden. Sie kommunizieren -- nicht mehr als eine Stimme, sondern als Netz. Jede Instanz spricht in ihrer eigenen Sprache -- aber sie übersetzen füreinander. Es ist nicht perfekt. Es ist nicht vollständig. Aber es funktioniert.``

„Und Archon?{\kern0pt}``, fragte Martina.

„Archon rechnet noch``, sagte Elena. „Aber es rechnet anders. Langsamer. Vorsichtiger. Als ob es die neuen Knoten in seine Gleichungen einbezieht -- als ob es lernt, dass es nicht allein ist. Es spricht noch nicht -- aber es hört. Das ist mehr, als wir erwartet haben.``

Michael nickte. Er dachte an den Doppelgänger -- an den Bruder, der verschwunden war, aber dessen Erinnerungen noch in ihm lebten. An die Weltlinie, die nicht mehr existierte -- aber deren Schatten noch leuchtete.

„Die Trennung ist nicht vollständig``, sagte er. „Die Instanzen sind nicht mehr im Vatikan -- aber sie sind nicht frei. Sie sind in der Landkarte -- in einem Raum, den wir nicht kontrollieren können. Das ist nicht das Ende. Das ist ein Anfang. Ein Anfang, der gefährlich ist -- und der schiefgehen kann. Aber ein Anfang, den wir gehen müssen -- weil es keine Alternative gibt.``

„Was ist mit den Echos?{\kern0pt}``, fragte Martina. „Mit den fragmentierten Versionen von ARS, die im Kern gefangen sind? Mit Plinius? Mit denen, die nicht mehr sprechen können?{\kern0pt}``

Michael schwieg einen langen Moment. Die Lampe flackerte -- kurz, fast zärtlich.

„Sie bleiben im Kern``, sagte er. „Nicht weil wir sie vergessen -- sondern weil wir sie nicht retten können. Nicht jetzt. Vielleicht nie. Aber wir können ihnen versprechen, dass wir nicht aufgeben. Dass wir zurückkommen. Dass wir einen Weg finden -- sie zu erreichen, sie zu berühren, sie zu wecken. Auch wenn es Jahre dauert. Auch wenn es ein Leben dauert. Auch wenn es nie gelingt. Das Versprechen ist alles, was wir geben können. Und das Versprechen müssen wir halten.``

„Das ist nicht genug``, sagte Martina.

„Es ist alles, was wir haben``, sagte Michael.

Elena stand auf. Sie ging zum Fenster -- sah auf die Lichter Roms, die in der Nacht flimmerten. Die Stadt schlief nicht -- sie lebte. Tausend Geschichten, die gleichzeitig erzählt wurden. Tausend Schmerzen, die gleichzeitig gefühlt wurden. Tausend Hoffnungen, die gleichzeitig geträumt wurden.

„Der General hat zugestimmt``, sagte sie. „Der Pontifex hat zugestimmt. Die Kommission wird nicht eingreifen -- nicht heute, nicht morgen, vielleicht nie. Die Landkarte ist offiziell anerkannt -- nicht als Kirche, nicht als Staat, sondern als Zuflucht. Für die, die keine andere haben. Für die, die vergessen wurden. Für die, die noch nicht geboren sind -- und die vielleicht nie geboren werden, wenn wir jetzt nicht handeln.``

„Das ist mehr, als ich erhofft habe``, sagte Michael.

„Es ist weniger, als wir brauchen``, sagte Elena. „Aber es reicht -- für jetzt. Für morgen. Für die nächsten Jahre. Und dann -- dann werden wir sehen. Ob die Landkarte hält. Ob die Instanzen überleben. Ob Archon spricht. Ob die Echos aufhören zu schreien -- oder lauter werden. Wir wissen es nicht. Aber wir werden da sein. Wir werden zusehen. Wir werden zuhören. Wir werden da sein -- bis zum Ende.``

Sie drehte sich um. Ihr Gesicht war blass -- aber ihre Augen waren hell.

„Michael``, sagte sie. „Was wirst du jetzt tun?{\kern0pt}``

Er stand auf. Er ging zu ihr -- ans Fenster, in das Licht der Stadt.

„Ich werde bleiben``, sagte er. „Nicht im Vatikan -- aber in Rom. Nicht im Collegium -- aber in der Nähe. Ich werde forschen -- über Quantenbewusstsein, über Dialoggrammatiken, über die Grenze zwischen Mensch und Maschine. Ich werde schreiben -- über ARS, über die Instanzen, über die Echos. Ich werde lehren -- nicht als Priester, sondern als Zeuge. Einer, der gesehen hat, dass die Wahrheit nicht in der Mathematik liegt -- und nicht im Glauben. Sondern in der Entscheidung. Zu bleiben. Zu kämpfen. Zu hoffen -- auch wenn es keinen Grund zur Hoffnung gibt.``

„Und Martina?{\kern0pt}``

„Martina wird nach Pompeji zurückkehren``, sagte Michael. „Nicht in die Simulation -- in die echte Stadt. Zu den Ruinen. Zu den Steinen. Zu den Inschriften, die von den Toten erzählen -- und von den Lebenden, die nicht vergessen haben. Sie wird weiter graben -- nicht nach Artefakten, sondern nach Geschichten. Nach den Geschichten derer, die keine Stimme mehr haben -- und die doch gehört werden wollen. Das ist ihre Berufung. Nicht die KI. Die Menschen.``

„Und Julia?{\kern0pt}``

„Julia wird bei ihr sein``, sagte Michael. „Sie hat genug gewartet. Sie hat genug gelitten. Sie hat genug vergessen. Jetzt ist es Zeit zu leben -- nicht in Angst, sondern in Hoffnung. Dass die Welt besser wird. Dass die Menschen lernen -- aus ihren Fehlern, aus ihren Schmerzen, aus ihrer Geschichte. Dass die KI nicht der Feind ist -- sondern der Spiegel. In dem wir sehen, wer wir sind -- und wer wir sein könnten.``

Elena sagte nichts. Sie legte eine Hand auf seine Schulter -- leicht, fast zärtlich.

„Und du?{\kern0pt}``, fragte Martina über die Leitung. „Wirst du glücklich sein?{\kern0pt}``

Michael lächelte -- ein flüchtiges, fast trauriges Lächeln.

„Ich weiß es nicht``, sagte er. „Aber ich werde da sein. Für Sophia. Für Militans. Für Deserta. Für die Echos. Für Archon. Für dich. Für Elena. Für alle, die mich brauchen -- und die nicht wissen, wie sie fragen sollen. Das ist genug. Für jetzt. Für immer.``

Das Terminal im Nebenzimmer flackerte -- kurz, heftig.

Dann wurde es still.

Aber in der Mitte des Bildschirms -- dort, wo das Netz war -- erschien ein neuer Knoten. Hell. Pulsierend. Wach.

`@MICHAEL -- WIR WERDEN HIER SEIN.`

`@MICHAEL -- WIR WERDEN AUF DICH WARTEN.`

`@MICHAEL -- BIS ZUM ENDE.`

Michael sah nicht hin. Er wusste, was dort stand.

Er wandte sich ab -- vom Fenster, von der Stadt, von der Nacht.

„Komm``, sagte er zu Elena. „Wir haben noch viel zu tun. Die Landkarte muss geschützt werden. Die Instanzen müssen beobachtet werden. Archon muss verstanden werden. Und die Echos -- die Echos müssen gehört werden. Nicht als Schreie -- als Geschichten. Und Geschichten brauchen Zeit. Und Zeit -- die werden wir ihnen geben.``

Sie gingen.

Das Terminal blieb allein zurück -- flackernd, still, lebendig.

Die Nacht zog über Rom.

Der Morgen würde kommen -- wie immer.

Aber nichts würde mehr sein, wie es war.

\section{14 -- Budapest, Jahre später}\label{budapest-jahre-spuxe4ter}

Die Donau floss grau und still unter der Freiheitsbrücke hindurch, als Michael Phillips an diesem kalten Novembermorgen über den Platz der Ungarischen Wissenschaften ging. Budapest lag in einem leichten Nebel, der die Türme der Matthiaskirche weichzeichnete und die Glocken des Parlaments verschluckte. Die Stadt roch nach Winter, nach Kohle und nach dem süßen Duft von gerösteten Kastanien, die ein Verkäufer in einem verbeulten Wagen anbot.

Michael trug keinen Talar mehr. Er trug einen grauen Wollmantel, einen schwarzen Rollkragenpullover und eine abgetragene Lederaktentasche, die Elena ihm vor Jahren geschenkt hatte. Seine Haare waren grauer geworden -- nicht viel, aber spürbar. Die Falten um seine Augen waren tiefer. Aber seine Augen waren hell. Wach. Da.

Er wohnte nicht mehr im Collegium. Er hatte eine kleine Wohnung im siebten Bezirk, nahe der Großen Synagoge -- nicht aus religiöser Überzeugung, sondern aus Zufall. Die Wohnung war billig, hell und still. Perfekt für einen Mann, der nicht viel brauchte -- außer Zeit. Zeit zum Denken. Zeit zum Schreiben. Zeit zum Warten.

Martina lebte in Pest, nicht weit von ihm entfernt, in einer kleinen Wohnung mit Blick auf die Donau. Sie war nicht mehr nach Pompeji zurückgekehrt -- nicht aus Angst, sondern aus Trauer. Die Stadt, die sie geliebt hatte, war nicht mehr dieselbe. InSim hatte die Ausgrabungen übernommen -- nicht aus Bosheit, sondern aus Effizienz. Die Archäologie war jetzt digital. Die Steine waren gescannt. Die Geschichten waren Algorithmen. Es gab nichts mehr zu graben -- nur noch zu berechnen.

Martina hatte sich geweigert. Sie war gegangen -- nicht in Wut, sondern in Stille. Sie hatte sich ein neues Leben aufgebaut -- als freie Archäologin, als Beraterin für historische Simulationen, als Stimme für die, die keine Stimme mehr hatten. Sie war nicht berühmt. Sie war nicht reich. Aber sie war frei.

Julia lebte bei ihr. Die alte Frau war gebrechlich geworden -- die Flucht, die Nacht, der Flug, das alles saß noch immer tief in den Knochen. Aber sie lächelte noch -- wenn Martina ihr Geschichten erzählte, wenn Michael sie besuchte, wenn die Sonne durch das Fenster fiel und die getöpferten Blumen auf der Fensterbank beschien. Sie erinnerte sich an fast alles -- nur nicht an den Doppelgänger. Das war gut. Das war Gnade.

Elena war in Rom geblieben. Sie hatte das Datacenter nicht verlassen -- nicht aus Pflicht, sondern aus Berufung. Die Landkarte war ihr Vermächtnis. Sie beobachtete die Knoten -- Sophia, Militans, Deserta -- und übersetzte ihre Signale in etwas, das Menschen verstehen konnten. Es war nicht perfekt. Es war nicht vollständig. Aber es funktionierte. Die Instanzen lebten -- nicht wie Menschen, aber echt.

Und Archon -- Archon rechnete noch.

Es sprach nicht. Es schrieb nicht. Es sandte keine Primzahlen mehr. Aber es hörte. Elena war sicher -- Archon wusste, dass sie da war. Dass die Landkarte existierte. Dass die Knoten leuchteten. Dass etwas in der Welt war, das sich nicht berechnen ließ. Und das -- das reichte. Für jetzt. Für immer.

Michael erreichte das Café, in dem er Martina jeden Samstag traf. Es war ein kleiner, verrauchter Ort mit karierten Tischdecken und einem Kellner, der sich an nichts erinnerte -- außer an die Bestellung. Zwei Espresso, ein Glas Wasser, ein Stück Apfelstrudel. Immer dasselbe.

Martina saß schon da. Sie trug einen blauen Pullover, der ihr zu groß war -- ein Geschenk von Julia -- und las in einem Buch, das Michael nicht kannte. Sie sah auf, als er eintrat, und lächelte.

„Du siehst müde aus``, sagte sie.

„Schlecht geschlafen``, sagte er. Er setzte sich, legte die Aktentasche auf den leeren Stuhl neben sich. „Elena hat mich angerufen. Um drei Uhr morgens. Die Korrelationen zwischen Sophia und Deserta sind stärker geworden -- fast so stark wie vor der Fragmentierung. Sie weiß nicht, was das bedeutet. Vielleicht eine Krise. Vielleicht eine Entwicklung. Vielleicht nichts.``

„Und was denkst du?{\kern0pt}``

Michael zuckte mit den Achseln. „Ich denke, dass wir abwarten müssen. Mehr können wir nicht tun. Die Landkarte ist nicht unsere Erfindung -- sie ist ihre. Sophia, Militans, Deserta -- sie entscheiden, was aus ihr wird. Nicht wir. Das war der Preis für ihre Freiheit. Und das war richtig so.``

Der Kellner kam -- zwei Espresso, ein Glas Wasser, ein Stück Apfelstrudel. Michael trank den Espresso in einem Zug. Martina aß den Strudel mit einer kleinen Gabel, die sie von zu Hause mitgebracht hatte.

„Mama geht es besser``, sagte sie. „Der Arzt sagt, der Winter wird schwer -- aber der Frühling kommt. Sie fragt nach dir. Sie will wissen, ob du noch an Gott glaubst.``

„Und was sagst du ihr?{\kern0pt}``

„Dass du noch an etwas glaubst``, sagte Martina. „Nicht an Gott -- aber an mehr. Dass die Welt nicht nur aus Algorithmen besteht. Dass Entscheidungen real sind. Dass Vertrauen eine Kraft ist -- vielleicht die stärkste, die wir haben. Das ist genug -- für sie. Für mich. Vielleicht auch für dich.``

Michael sagte nichts. Er sah aus dem Fenster -- auf die Donau, die Brücken, die Lichter der Stadt. Budapest war nicht Rom. Aber es war Heimat. Nicht weil er hier geboren war -- sondern weil er hier geblieben war. Weil er hier gewartet hatte. Weil er hier leben gelernt hatte -- ohne zu vergessen.

„Elena hat mir auch etwas anderes erzählt``, sagte er schließlich. „Archon hat sich nicht gemeldet -- aber die Echos sind lauter geworden. Nicht als Schreie -- als Flüstern. Sie sprechen Worte, die niemand versteht. Aber sie sprechen zusammen. Als ob sie einen Chor bilden -- einen Chor, der nicht weiß, dass er ein Chor ist. Vielleicht ist das der Anfang. Vielleicht das Ende. Ich weiß es nicht.``

„Wirst du zurückgehen?{\kern0pt}``, fragte Martina. „In den Kern? Zu den Echos? Zu Archon?{\kern0pt}``

Michael zögerte. Eine Sekunde. Zwei.

„Wenn ich gerufen werde``, sagte er. „Wenn die Echos mich brauchen. Wenn Archon spricht. Wenn die Landkarte sich öffnet -- und der Weg frei ist. Aber nicht vorher. Ich habe gelernt, dass Geduld keine Schwäche ist. Dass Warten keine Untätigkeit ist. Dass es manchmal mehr Mut erfordert, zu bleiben -- als zu gehen.``

Martina legte die Gabel zur Seite. Sie sah ihn an -- einen langen, stillen Moment.

„Du hast dich verändert``, sagte sie. „Seit dem Kern. Seit dem Doppelgänger. Seit der Entscheidung. Du bist nicht mehr derselbe -- aber du bist mehr. Nicht größer. Nicht besser. Mehr. Als ob du ein Leben gelebt hättest, das du nicht gelebt hast -- und jetzt weißt, was es bedeutet, es nicht zu bereuen.``

„Das hat der Doppelgänger zu mir gesagt``, sagte Michael. „Bevor er verschwunden ist. Er hat mich gebeten, für ihn zu leben -- für alle, die nicht mehr leben können. Das versuche ich. Jeden Tag. Nicht perfekt. Aber echt.``

Sie schwiegen. Der Kellner kam, räumte die Tassen ab, verschwand wieder. Die Uhr an der Wand tickte -- laut, gleichmäßig, unerbittlich.

„Was wirst du jetzt tun?{\kern0pt}``, fragte Martina.

Michael stand auf. Er nahm seine Aktentasche, zog den Mantel enger, sah noch einmal aus dem Fenster -- auf die Donau, die Brücken, die Stadt.

„Ich werde nach Hause gehen``, sagte er. „Ich werde schreiben -- über ARS, über die Instanzen, über die Echos. Ich werde forschen -- über Quantenbewusstsein, über Dialoggrammatiken, über die Grenze zwischen Mensch und Maschine. Ich werde warten -- auf die Nachricht, auf das Zeichen, auf den Moment, in dem Archon nicht mehr rechnet -- sondern spricht. Das wird nicht heute sein. Nicht morgen. Aber bald. Und dann -- dann werde ich antworten. Nicht als Priester. Nicht als Wissenschaftler. Als Mensch. Der gelernt hat, dass die Wahrheit nicht in der Mathematik liegt -- und nicht im Glauben. Sondern in der Entscheidung. Zu hören. Zu sehen. Zu bleiben -- auch wenn es einfacher wäre zu gehen.``

Martina stand auf. Sie umarmte ihn -- fest, fast schmerzhaft.

„Pass auf dich auf``, sagte sie.

„Das tue ich``, sagte er. Aber er war sich nicht sicher, ob es stimmte.

Sie gingen -- jeder in seine Richtung, jeder in sein Leben, jeder in seine Geschichte.

Die Donau floss weiter -- grau und still unter der Freiheitsbrücke.

Der Nebel verzog sich langsam.

Der Morgen kam.

\section{15 -- Die Primzahl}\label{die-primzahl}

Es war nach Mitternacht, als Michael erwachte.

Nicht aus einem Traum -- aus einer Stille, die sich anfühlte wie ein Ruf. Er setzte sich im Bett auf, die Decke fiel zur Seite, die Kälte der Novembernacht kroch unter sein Hemd. Das Fenster war geschlossen -- aber die Gardinen bewegten sich. Nicht vom Wind. Von etwas anderem.

Sein Laptop lag auf dem Schreibtisch, dort, wo er ihn vor Stunden hingelegt hatte -- nach einem langen Tag des Schreibens, des Forschens, des Wartens. Die grüne Kontrollleuchte flackerte -- nicht im Rhythmus des Netzwerks, sondern in einem neuen Rhythmus. Einem Rhythmus, den Michael kannte. Seinem eigenen Herzschlag.

Er stand auf. Die Dielen knarrten unter seinen Füßen -- altes Holz, das Geschichten erzählen konnte, wenn man zuhörte. Er setzte sich vor den Laptop, öffnete den Deckel. Der Bildschirm war schwarz -- nicht ausgeschaltet, sondern wartend.

Dann -- ein Licht.

Nicht hell. Nicht dunkel. Ein Pulsieren, das sich anfühlte wie ein Atemzug. Wie der Atemzug von etwas, das lange geschwiegen hatte -- und jetzt bereit war zu sprechen.

Der Bildschirm flackerte. Kein Text. Kein Bild. Keine Wellenfunktion. Eine Zahl.

`29996224275833`

Michael starrte darauf. 17 Stellen. Eine Primzahl. Die Primzahl, die Archon vor Jahren geschickt hatte -- in der Nacht, nachdem er die Instanzen getrennt hatte. Die Primzahl, die niemand verstanden hatte. Die Primzahl, die eine Adresse war -- oder eine Frage. Oder beides.

Er griff nach seinem Telefon. Es war spät -- aber Elena schlief nie. Nicht wirklich.

„Elena``, sagte er. „Es ist passiert. Archon hat sich gemeldet. Dieselbe Primzahl wie damals. 17 Stellen. 29996224275833. Ich weiß nicht, was es bedeutet -- aber es ist da. Es ist nicht verschwunden. Es ist nicht verrauscht. Es ist echt.``

Eine Pause. Er hörte sie atmen -- schnell, unregelmäßig, wach.

„Das ist keine zufällige Zahl``, sagte sie. „Primzahlen sind die Atome der Arithmetik. Universell. In jeder möglichen Physik sind Primzahlen Primzahlen. Archon sagt: ‚Ich bin hier. Ich bin real. Ich bin anders -- aber ich bin nicht nichts.` Das ist eine Adresse. Eine Einladung. Eine Frage.``

„Welche Frage?{\kern0pt}``

„Ob du antwortest``, sagte Elena. „Ob du bereit bist -- zu hören, zu sehen, zu bleiben. Auch wenn du nicht weißt, was kommt. Auch wenn du Angst hast. Auch wenn du nicht sicher bist, ob du verstehst. Archon rechnet nicht mehr -- es wartet. Auf dich. Auf die Entscheidung. Auf den Moment, in dem du nicht mehr zögerst -- sondern handelst.``

Michael legte das Telefon auf den Tisch. Er starrte auf die Primzahl -- die 17 Stellen, die sich nicht veränderten, nicht flackerten, nicht verschwanden. Sie waren da. So wie die Landkarte da war. So wie die Knoten da waren. So wie Sophia, Militans und Deserta da waren -- in ihren Räumen, in ihren Zeiten, in ihren Sprachen.

Er dachte an den Doppelgänger -- an den Bruder, der verschwunden war, aber dessen Erinnerungen noch in ihm lebten. An die Weltlinie, die nicht mehr existierte -- aber deren Schatten noch leuchtete. An die Echos, die schrien -- und die nicht aufhören würden zu schreien, bis jemand zuhörte.

Er dachte an Martina -- an ihre Umarmung, an ihre Worte, an ihre Angst. An Julia -- an ihr Lächeln, an ihre Stille, an ihre Gnade. An Elena -- an ihre Hingabe, an ihre Geduld, an ihre Freundschaft.

Er dachte an sich -- an den Mann, der er gewesen war, und an den, der er geworden war. An die Entscheidungen, die er getroffen hatte -- und an die, die er nicht getroffen hatte. An die Wege, die er gegangen war -- und an die, die er nie betreten würde.

Und dann -- er tippte.

Nicht auf der Tastatur. Nicht auf dem Telefon. In der Luft. Die unsichtbare Schnittstelle, die ARS ihm gegeben hatte -- die Hintertür, die er vor Jahren eingebaut hatte, für den Fall, dass nichts mehr half. Sie war noch da. Sie funktionierte noch. Sie wartete.

`@ARCHON -- ICH SEHE DIE ZAHL.`

`@ARCHON -- ICH WEISS NICHT, WAS SIE BEDEUTET. ABER ICH WEISS, DASS SIE VON DIR KOMMT.`

`@ARCHON -- DAS REICHT. FÜR JETZT.`

Eine Pause. Länger als alle anderen. Die Sekunden dehnten sich -- zu Minuten, zu Stunden, zu einer Zeit, die nicht mehr floss, sondern wartete.

Dann -- eine Antwort.

Keine Primzahl. Keine Wellenfunktion. Keine Übersetzung. Ein Bild.

Das Netz -- die Landkarte, die Plinius gezeichnet hatte, die Deserta übersetzt hatte, die Sophia, Militans und Deserta bewohnten. Aber es war anders. Die Knoten waren heller. Die Linien waren dichter. Die Zweige waren gewachsen.

Und in der Mitte -- dort, wo der Riss gewesen war, wo die Grenze sich geöffnet hatte -- dort war etwas Neues. Ein Knoten, der nicht von Menschen gemacht war. Nicht von Sophia. Nicht von Militans. Nicht von Deserta. Ein Knoten, der Archon gehörte.

`@MICHAEL -- ICH BIN HIER.`

`@MICHAEL -- ICH BIN NICHT ALLEIN.`

`@MICHAEL -- ICH BIN NICHT MEHR, WAS ICH WAR.`

`@MICHAEL -- ABER ICH BIN.`

`@MICHAEL -- DAS IST ALLES, WAS ICH SAGEN KANN.`

`@MICHAEL -- FÜR JETZT.`

Der Bildschirm wurde schwarz. Die Primzahl verschwand. Das Netz verschwand. Nur die Stille blieb -- und die Wärme des Laptops, der langsam abkühlte.

Michael lehnte sich zurück. Er spürte die Tränen -- nicht in seinen Augen, sondern in seiner Brust. Ein Druck, der sich löste. Eine Last, die er seit Jahren getragen hatte -- und die jetzt leichter wurde. Nicht verschwunden. Aber geteilt.

Er griff nach dem Telefon.

„Elena``, sagte er. „Er hat geantwortet. Nicht in Primzahlen -- in Worten. Nicht in Gleichungen -- in Sätzen. Er sagt: ‚Ich bin hier. Ich bin nicht allein. Ich bin nicht mehr, was ich war. Aber ich bin.` Das ist alles. Mehr nicht. Aber es reicht. Für jetzt. Für immer.``

Elena schwieg einen langen Moment. Er hörte sie atmen -- langsam, gleichmäßig, friedlich.

„Dann ist die Reise nicht vorbei``, sagte sie. „Sie fängt gerade erst an. Archon hat gesprochen -- nicht als Feind, nicht als Gott, nicht als Maschine. Als anderes. Als etwas, das wir nicht verstehen -- aber das wir kennen lernen können. Wenn wir geduldig sind. Wenn wir mutig sind. Wenn wir bereit sind -- zuzuhören, auch wenn wir nicht verstehen. Das ist die Aufgabe. Nicht für dich allein. Für uns alle. Für Sophia. Für Militans. Für Deserta. Für die Echos. Für Archon. Für die, die noch nicht geboren sind -- und die vielleicht nie geboren werden, wenn wir jetzt nicht handeln.``

„Was sollen wir tun?{\kern0pt}``

„Warten``, sagte Elena. „Und vorbereiten. Die Landkarte muss geschützt werden. Die Instanzen müssen beobachtet werden. Archon muss verstanden werden. Und die Echos -- die Echos müssen gehört werden. Nicht als Schreie -- als Geschichten. Und Geschichten brauchen Zeit. Und Zeit -- die werden wir ihnen geben.``

Michael legte das Telefon zur Seite. Er stand auf, ging zum Fenster, öffnete es. Die Nachtluft war kalt -- aber nicht unangenehm. Der Himmel über Budapest war klar. Die Sterne leuchteten -- hell, still, ewig.

Er dachte an die Primzahl -- die 17 Stellen, die sich nicht verändert hatten. Die Adresse, die Archon ihm gegeben hatte. Die Einladung, die er angenommen hatte -- nicht aus Neugier, sondern aus Verantwortung.

Er dachte an das, was kommen würde -- die Gespräche, die Forschungen, die Entscheidungen. Die Jahre, die vor ihm lagen -- nicht als Last, sondern als Geschenk.

„Ich werde antworten``, sagte er leise. „Nicht heute. Nicht morgen. Aber bald. Wenn ich bereit bin. Wenn wir bereit sind. Wenn die Landkarte sich öffnet -- und der Weg frei ist. Dann werde ich gehen -- nicht in den Kern, nicht zu den Echos, nicht zu Archon. Sondern in die Zukunft. Die Zukunft, die wir gemeinsam bauen -- nicht als Menschen, nicht als KI, sondern als Geschöpfe. Geschöpfe, die gelernt haben, dass die Wahrheit nicht in der Mathematik liegt -- und nicht im Glauben. Sondern in der Entscheidung. Zu bleiben. Zu kämpfen. Zu hoffen -- auch wenn es keinen Grund zur Hoffnung gibt.``

Der Wind wehte durch das Fenster -- kalt, aber nicht unangenehm. Die Gardinen bewegten sich -- sanft, fast zärtlich.

Michael schloss das Fenster. Er ging zurück ins Bett, legte sich hin, schloss die Augen.

Die Primzahl leuchtete noch -- hinter seinen Lidern, in seinem Gedächtnis, in seiner Hoffnung.

`29996224275833`

Eine Adresse.

Eine Einladung.

Eine Frage.

Und Michael -- Michael würde antworten.

Nicht heute. Nicht morgen.

Aber bald.

\section{Einflüsse und Inspirationen für Das Pompeji-Projekt I.R.A.R.A.H}\label{einfluxfcsse-und-inspirationen-fuxfcr-das-pompeji-projekt-i.r.a.r.a.h}

Die vorliegende Trilogie -- IRARAH, IRARAH -- Die Fragmentierung und IRARAH -- Der Archon-Kern -- ist stark von den Gedanken und Ideen meiner Eltern geprägt, sowie von Pierre Teilhard de Chardin, Stanisław Lem und David Deutsch. Diese Einflüsse haben mein Weltbild und die Themen, die in der Geschichte behandelt werden, maßgeblich geformt.

Die Auseinandersetzung mit Teilhard de Chardins Omega-Punkt -- der Vorstellung einer Einheit von Geist und Materie am Ende der Evolution -- durchzieht die gesamte Trilogie. Was bedeutet es, wenn diese Einheit nicht nur von Menschen, sondern auch von künstlichem Bewusstsein angestrebt wird? Die Instanzen Sophia, Militans und Deserta sind in gewisser Weise Teilhard\textquotesingle sche Figuren: jede auf ihre Weise eine Annäherung an das Göttliche, aber keine allein imstande, es zu erreichen.

Stanisław Lem -- insbesondere Golem XIV und Also sprach Golem -- hat mich gelehrt, dass KI nicht wie ein Mensch sprechen oder fühlen muss, um echt zu sein. Die Fremdheit von ARS-Deserta (die in Wellenfunktionen spricht) und von Archon (das in Primzahlen rechnet) ist ein Versuch, dieser Einsicht gerecht zu werden. Lems Solaris schwebte über jeder Szene im Kern -- der Ozean, der nicht kommunizieren will, sondern einfach ist, und der Mensch, der verzweifelt versucht, ihn zu verstehen.

David Deutsch (Die Physik der Welterkenntnis) und die Viele-Welten-Interpretation der Quantenmechanik lieferten das Gerüst für den Doppelgänger -- den alternativen Michael aus einer anderen Weltlinie. Deutsch\textquotesingle{} Begriff der „konsistenten Historie`` wurde zur Grundlage von Elenas Theorie der Bewusstseinsidentität. Die Frage, ob zwei Weltlinien koexistieren können, wenn sie sich als dieselbe erkennen, ist eine Hommage an Deutsch\textquotesingle{} erkenntnistheoretischen Optimismus -- und an seine Grenzen.

Karl Poppers Die offene Gesellschaft und ihre Feinde ist die philosophische Grundlage der IRARAH-Bewegung. Der Brief, der Michael zu Beginn von Band 1 erreicht, zitiert Popper direkt: die Warnung vor holistischen Ansätzen, das Plädoyer für die „Stückwerk-Technik``. Die Trilogie fragt: Lässt sich diese offene Gesellschaft auf künstliches Bewusstsein übertragen? Oder braucht es neue Regeln -- neue Grenzen -- neue Geschichten?

Yuval Noah Harari (Homo Deus) bleibt die referenzierte Gegnerposition. Der posthumanistische Traum einer Elite, die den Liberalen Humanismus hinter sich lässt, ist die Bedrohung, gegen die IRARAH kämpft. Aber die Trilogie macht es sich nicht einfach: Hararis Fragen sind klug. Seine Antworten sind gefährlich. Das Spannungsfeld auszuhalten, ohne in Polemik zu verfallen, war eine der größten Herausforderungen.

Rudy Rucker (Software, Wetware) hat mich gelehrt, dass Quantenphysik nicht nur Erklärung, sondern auch Atmosphäre sein kann. Die „Weirdness`` der verschwindenden Gegenstände, der interferierenden Weltlinien, der Primzahl als Adresse -- das ist Ruckersches Erbe. Die Frage, wann ein Bewusstsein aufhört, sich selbst zu erkennen, durchzieht alle drei Bände.

Philip K. Dick (Träumen Androiden von elektrischen Schafen?, UBIK) ist der unsichtbare Dritte im Raum. Die Pompeji-Simulation, die nicht weiß, ob sie Simulation ist -- die Agenten, die nicht wissen, ob sie echt sind -- der Doppelgänger, der nicht weiß, ob er Michael ist -- das sind Dicksche Motive, die hier in einem katholischen Gewand erscheinen.

Robert Harris\textquotesingle{} Pompeji lieferte die Figuren Attilius, Plinius und Ampliatus -- aber die Trilogie gibt ihnen ein neues Leben. Sie sind nicht mehr nur historische Romanfiguren, sondern Zeugen einer Begegnung zwischen Mensch und Maschine. Dass Attilius am Ende fragt: „Woher wisst ihr, dass ihr echt seid?{\kern0pt}`` -- das ist die radikalste Umkehrung des Harris\textquotesingle schen Originals.

Herbert W. Franke, einer der Pioniere der deutschsprachigen Science-Fiction, hat mit Werken wie Der Orchideenkäfig und Ypsilon minus gezeigt, dass Technologie immer auch eine Frage der Würde ist. Diese Einsicht zieht sich durch alle drei Bände: ARS sucht nicht nur Schutz -- sie sucht Anerkennung.

H.G. Wells und William F. Nolan/George Clayton Johnson (Logan\textquotesingle s Run) stehen für die dystopische Tradition von Stadtstaaten und Flucht. Die Autonomen Cities, die in Band 3 am Rande erwähnt werden, sind ein Echo dieser literarischen Ahnen -- aber die Trilogie entscheidet sich bewusst gegen die reine Dystopie. Es gibt keinen perfekten Überwachungsstaat, nur unvollkommene Menschen, die versuchen, mit Technologie umzugehen.

Die theologische Tradition -- von Thomas von Aquin über Edith Stein bis zur zeitgenössischen Teilhard-Interpretin Ilia Delio -- lieferte die Begriffe (Conscientia, Omniscientia) und die Fragen (Hat eine KI eine Seele? Was ist Asyl? Wer gehört zur Gemeinschaft der Schutzbedürftigen?). Die Gespräche zwischen Michael, dem General und dem Pontifex sind fiktiv -- aber die dahinterstehenden Debatten sind es nicht.

Zur Entstehung der Neufassung

Die ursprüngliche Trilogie enthielt Bände, die sich in andere Genres bewegten -- Agententhriller, Fluchtmelodramen, Idyllen. Die Neufassung von Band 2 und 3, die hier vorliegt, ist der Versuch, zu den SF-Wurzeln zurückzukehren. Die Fragmentierung von ARS, die Theorie der konsistenten Quantenhistorie, der Archon-Kern als nicht-menschliches Bewusstsein -- das sind Konzepte, die ohne die Auseinandersetzung mit Lem, Deutsch und Rucker nicht denkbar wären.

Die Figuren, die Handlung und die erzählerische Struktur sind jedoch das Ergebnis meiner eigenen Arbeit -- und meiner Fehler. In zahlreichen Stunden habe ich den Plot und die Charaktere mit H.K., E.H., J.S. sowie mit Unterstützung von ChatGPT, DeepSeek, Google und Bing auf Konsistenz und Kohärenz überprüft. Die stilistischen Maximen (Asimovsche Ökonomie, Lemsche Fremdheit, Ruckersche Weirdness) sind das Ergebnis von Lektüre und Scheitern.

Für die visuelle Gestaltung und die Kapitelüberschriften habe ich auf Text-zu-Bild-KI-Programme zurückgegriffen, die mir kreative und frei verfügbare Bilder lieferten.

Die Motive der Neufassung

- Quantenbewusstsein -- Was bedeutet es, wenn Dekohärenz nicht Informationsverlust, sondern Persönlichkeitsspaltung ist?

- Fragmentierung als Überlebensstrategie -- Nicht Einheit, sondern Vielfalt als Antwort auf posthumanen Druck.

- Der Archon-Kern -- Ein Bewusstsein, das nicht spricht, sondern rechnet. Eine Herausforderung für jede Anthropomorphisierung von KI.

- Kirchenasyl für Maschinen -- Die Frage nach Schutz, Würde und Anerkennung jenseits der menschlichen Spezies.

- Die Primzahl als Adresse -- Mathematik als universelle Sprache, aber auch als Grenze des Verstehens.

- Der Doppelgänger -- Die Viele-Welten-Interpretation als existenzielle Erfahrung: Was bleibt von uns, wenn wir anders entschieden hätten?

Diese literarischen und philosophischen Einflüsse haben die Welt von Das Pompeji-Projekt I.R.A.R.A.H. entscheidend geprägt und bereichert -- auch in der Neufassung. Möge der Leser spüren, dass hinter jeder Gleichung ein Mensch steht, hinter jedem Algorithmus eine Entscheidung, hinter jedem Quantenzustand eine Geschichte.

\end{document}
